% FILE: Registry of ongoing and planned ME/CFS research studies
\chapter{Ongoing and Planned ME/CFS Research Studies}
\label{app:research-registry}

This appendix catalogues major ongoing and planned research studies in Myalgic
Encephalomyelitis/Chronic Fatigue Syndrome (ME/CFS).  Each entry provides
structured metadata---principal investigator, institution, registry
identifiers, cohort size, and study focus---followed by a brief discussion of
the study's potential to advance ME/CFS understanding.

\textbf{Data sources.}  Entries were compiled from ClinicalTrials.gov, the
ISRCTN registry, institutional announcements, and the Solve ME/CFS Initiative
trial tracker.  Registry identifiers (NCT numbers, ISRCTN numbers) are
included where available.

\textbf{Currency.}  This registry was last updated on 2026-03-07.
Study status information may have changed since that date; readers should
verify current status via the URLs provided.

\textbf{Scope.}  The registry includes studies that are actively recruiting,
enrolled, or in analysis as of the currency date, plus selected planned studies
with confirmed funding.  Completed studies with published results are covered
in the main text and Appendix~\ref{app:annotated-bibliography}.

\begin{limitation}[Registry Completeness]
This registry is necessarily incomplete.  Smaller institutional studies,
industry-sponsored trials in early phases, and studies registered outside
ClinicalTrials.gov or ISRCTN may not be captured.  Inclusion does not imply
endorsement of study design or methodology; exclusion does not imply
irrelevance.
\end{limitation}

% =============================================================================
\section{Multi-Omics and Systems Biology}
\label{sec:registry-multiomics}
% =============================================================================

\subsection{Rosetta Stone Study}
\label{sec:registry-rosetta-stone}

\begin{description}
    \item[Principal Investigator:] Danny Altmann, Professor of Immunology
    \item[Institution:] Imperial College London, United Kingdom
    \item[Contact/URL:] \url{https://meassociation.org.uk/2025/12/driving-discovery-the-me-association-invests-1-1m-into-pioneering-research-programme/}
    \item[Funder:] ME Association, Ramsay Research Fund
    \item[Status:] Active (recruitment began late 2025)
    \item[Phase:] Observational, multi-omics
    \item[Cohort:] 250 ME/CFS + 250 Long Covid + matched healthy controls
    \item[Mechanism/Focus:] Side-by-side cellular and molecular analyses of
        ME/CFS and Long Covid.  Stool, blood, and saliva samples analysed
        across multiple omics platforms, with machine learning applied to
        discover shared molecular pathways between the two conditions.
        Partnership with ELAROS for digital phenotyping.
    \item[Primary Outcomes:] Shared and distinguishing molecular signatures
        between ME/CFS and Long Covid; candidate biomarkers
    \item[Estimated Completion:] 2028 (3-year study)
    \item[Timeline:] Recruitment started Q4 2025; 3-month progress update
        published March 2026; interim analysis expected mid-study
    \item[Publication Medium:] Not yet announced; likely high-impact
        immunology or multidisciplinary journal; open access expected given
        charity funding
    \item[Document Relevance:] Ch.~7 (immune dysfunction), Ch.~11 (gut
        microbiome), Ch.~13 (integrative models), Ch.~14 \S\ref{sec:cross-disease}
        (cross-disease comparison)
\end{description}

The Rosetta Stone Study represents the largest single investment in ME/CFS
biomedical research by a UK charity (\pounds1.1M).  Its defining strength is
the parallel recruitment of ME/CFS and Long Covid cohorts with identical
multi-omics profiling, enabling direct molecular comparison of two conditions
widely hypothesised to share pathophysiological mechanisms.  If the study
identifies shared immune or metabolic signatures, it could accelerate
biomarker development by leveraging the larger Long Covid research
infrastructure.  The inclusion of stool samples positions the study to
contribute microbiome data that remains scarce in well-phenotyped ME/CFS
cohorts.

\subsection{BioQuest Study}
\label{sec:registry-bioquest}

\begin{description}
    \item[Principal Investigator:] Jonas Bergquist (proteomics lead);
        coordinated by Open Medicine Foundation (OMF) research team
    \item[Institution:] Uppsala University (Sweden) and Open Medicine
        Foundation collaborative centres
    \item[Contact/URL:] \url{https://ns1.omf.ngo/bioquest-update/}
    \item[Funder:] Open Medicine Foundation
    \item[Status:] Active (sample collection complete; analysis phase)
    \item[Phase:] Observational, multi-omics biomarker discovery
    \item[Cohort:] 400 ME/CFS + 400 healthy controls + 200 disease controls
        (50 each: multiple sclerosis, burnout syndrome, depression, Long Covid).
        Combines samples from the Uppsala Collaborative Centre and the NIH
        Chronic Fatigue Initiative biobank.
    \item[Mechanism/Focus:] Large-scale metabolomics, lipidomics, proteomics,
        and cytokine profiling (\textgreater10,000 analytes) with AI-driven
        subgroup identification.  Goal: identify a 5--20 protein/metabolite
        biomarker panel unique to ME/CFS.
    \item[Primary Outcomes:] Validated ME/CFS biomarker panel; biologically
        defined patient subgroups
    \item[Estimated Completion:] Results expected 2026--2027
    \item[Timeline:] Sample integration approved 2025; testing began late 2025;
        bioinformatic analysis ongoing
    \item[Publication Medium:] Expected in high-impact journal; OMF typically
        publishes open access
    \item[Document Relevance:] Ch.~6 (energy metabolism), Ch.~7 (immune
        dysfunction), Ch.~13 (integrative models), Ch.~20 (biomarker research)
\end{description}

BioQuest is the largest ME/CFS biomarker study to date by sample size
($n=1{,}000$).  Its inclusion of disease controls (MS, depression, burnout,
Long Covid) directly addresses a longstanding criticism that ME/CFS biomarker
candidates lack disease specificity.  The multi-omics approach measuring over
10,000 analytes, combined with AI-driven clustering, has realistic potential to
identify biologically defined subgroups---a prerequisite for precision
medicine approaches.  If successful, BioQuest could provide the first
blood-based diagnostic panel capable of distinguishing ME/CFS from conditions
with overlapping symptom profiles.

% =============================================================================
\section{Genomics and Epigenetics}
\label{sec:registry-genomics}
% =============================================================================

\subsection{DecodeME}
\label{sec:registry-decodeme}

\begin{description}
    \item[Principal Investigator:] Chris Ponting, Professor of Medical
        Bioinformatics
    \item[Institution:] University of Edinburgh, United Kingdom
    \item[Contact/URL:] \url{https://www.decodeme.org.uk/}
    \item[Registry ID:] ISRCTN15960771
    \item[Funder:] MRC/NIHR (UK Research and Innovation)
    \item[Status:] Active (genotyping and analysis phase)
    \item[Phase:] Observational, genome-wide association study (GWAS)
    \item[Cohort:] \textgreater20,000 ME/CFS patients (DNA samples);
        UK-based, self-reported diagnosis with GP confirmation
    \item[Mechanism/Focus:] First large-scale GWAS for ME/CFS.  Aims to
        identify genetic variants associated with susceptibility, severity,
        and symptom subtypes.  DNA collected via saliva kits mailed to
        participants.
    \item[Primary Outcomes:] Genome-wide significant loci for ME/CFS;
        genetic correlation with related conditions; polygenic risk scores
    \item[Estimated Completion:] First results expected 2026
    \item[Timeline:] Recruitment launched 2022; sample collection completed
        2024; genotyping in progress; first GWAS results anticipated 2026
    \item[Publication Medium:] Expected in genetics/genomics journal
        (e.g., \textit{Nature Genetics}); open access
    \item[Document Relevance:] Ch.~12 (genetic and epigenetic factors),
        Ch.~20 (biomarker research), Ch.~23 (epidemiology)
\end{description}

DecodeME is unprecedented in ME/CFS genetics: with over 20,000 participants,
it has the statistical power to detect common variants of modest effect size
that previous underpowered GWAS could not identify.  Genetic findings could
reveal causal pathways (e.g., immune regulation, ion channel function,
autonomic control) and enable Mendelian randomisation analyses to test whether
associations observed in smaller studies reflect causal relationships.
The study's scale also allows subgroup analyses by onset trigger, severity,
and symptom profile.

\subsection{EpiSwitch CFS Epigenetic Diagnostic}
\label{sec:registry-episwitch}

\begin{description}
    \item[Principal Investigator:] Ewan Hunter (lead author, validation study)
    \item[Institution:] Oxford BioDynamics, United Kingdom
    \item[Contact/URL:] \url{https://doi.org/10.1186/s12967-025-07203-w}
    \item[Registry ID:] N/A (proof-of-concept; clinical validation pending)
    \item[Funder:] Oxford BioDynamics (industry)
    \item[Status:] Proof-of-concept published October 2025; clinical
        validation studies planned
    \item[Phase:] Diagnostic development (pre-clinical validation)
    \item[Cohort:] Initial study: 47 severe ME/CFS patients + 61 healthy
        controls.  Larger validation cohorts needed.
    \item[Mechanism/Focus:] EpiSwitch 3D genomic profiling of chromosome
        conformation captures (CCs) in blood.  Identified a 200-marker
        epigenetic model.  Pathways implicated: interleukins, TNF$\alpha$,
        neuroinflammation, toll-like receptor signalling, JAK/STAT.
    \item[Primary Outcomes:] Sensitivity 92\%, specificity 98\%, overall
        accuracy 96\% in proof-of-concept validation cohort
    \item[Estimated Completion:] Clinical validation timeline not yet
        announced
    \item[Timeline:] Proof-of-concept published October 2025; independent
        replication and prospective validation required before clinical use
    \item[Publication Medium:] \textit{Journal of Translational Medicine}
        (published); future validation studies TBD
    \item[Document Relevance:] Ch.~12 (genetic and epigenetic factors),
        Ch.~20 (biomarker research), Ch.~25 (translational findings)
\end{description}

The EpiSwitch proof-of-concept achieved striking diagnostic accuracy (96\%),
but the small sample size ($n=108$) and retrospective design warrant caution.
If validated in larger, prospective, multi-site cohorts with disease controls,
this approach could provide the first epigenetic blood test for ME/CFS.  The
implicated pathways (JAK/STAT, TLR signalling, neuroinflammation) align with
immune findings reported in Chapters~7 and~8, providing convergent evidence
from an independent methodological angle.  The critical next step is
independent replication with blinded samples.

\subsection{PrecisionLife Combinatorial Genomics}
\label{sec:registry-precisionlife}

\begin{description}
    \item[Principal Investigator:] Steve Gardner (CEO, PrecisionLife)
    \item[Institution:] PrecisionLife Ltd, Oxford, United Kingdom
    \item[Contact/URL:] \url{https://precisionlife.com/news-and-events/me-genetics-study}
    \item[Registry ID:] N/A (secondary analysis of DecodeME and UK Biobank data)
    \item[Funder:] PrecisionLife (industry) with access to DecodeME consortium
        data
    \item[Status:] Published December 2025
    \item[Phase:] Computational genomics (AI-led combinatorial analytics)
    \item[Cohort:] Two DecodeME cohorts + UK Biobank; results confirmed across
        three independent datasets
    \item[Mechanism/Focus:] Applied combinatorial analytics platform to
        identify gene--gene interaction networks in ME/CFS.  Unlike standard
        GWAS (which tests single variants), this approach identifies
        combinations of variants that jointly confer risk.
    \item[Primary Outcomes:] Identified 250+ core genes associated with
        ME/CFS, of which 76 genes overlap with long COVID.  Networks cluster
        around immune regulation, metabolic function, and neurological pathways.
    \item[Estimated Completion:] Published
    \item[Publication Medium:] PrecisionLife press release and preprint;
        peer-reviewed publication expected
    \item[Document Relevance:] Ch.~12 (genetic factors), Ch.~13 (integrative
        models), Ch.~14 \S\ref{sec:cross-disease} (long COVID overlap)
\end{description}

The PrecisionLife analysis complements the DecodeME GWAS
(Section~\ref{sec:registry-decodeme}) by detecting higher-order gene--gene
interactions that conventional single-variant approaches miss.  The 76-gene
overlap with long COVID provides the strongest genetic evidence yet for shared
biological mechanisms between the two conditions, supporting the hypothesis
explored in Ch.~14 \S\ref{sec:cross-disease}.  However, combinatorial
approaches can generate false positives when the number of tested combinations
is large; independent replication (ideally in non-European ancestries) is
essential before clinical or mechanistic conclusions are drawn.
\cite{PrecisionLife2025decodeme}

% =============================================================================
\section{Biomarker Discovery and Diagnostics}
\label{sec:registry-biomarker}
% =============================================================================

\subsection{BioMapAI (Gut Microbiome and Metabolomics)}
\label{sec:registry-biomapai}

\begin{description}
    \item[Principal Investigator:] Derya Unutmaz (Jackson Laboratory);
        AI platform developed with Duke University
    \item[Institution:] Jackson Laboratory (JAX) and Duke University
        School of Medicine, United States
    \item[Contact/URL:] \url{https://medschool.duke.edu/news/ai-thats-finally-making-sense-chronic-fatigue-syndrome}
    \item[Funder:] NIH and institutional funding
    \item[Status:] Active (longitudinal follow-up ongoing)
    \item[Phase:] Observational, AI-integrated multi-omics
    \item[Cohort:] 153 ME/CFS patients + 96 healthy controls, followed
        longitudinally over 4 years
    \item[Mechanism/Focus:] AI platform integrating gut metagenomics, plasma
        metabolomics, immune cell profiles, standard blood tests, and clinical
        symptoms.  Achieved 90\% accuracy in distinguishing ME/CFS from
        controls.  Identified gut microbiome signatures as strong predictors.
    \item[Primary Outcomes:] Multi-modal diagnostic classifier; longitudinal
        trajectory modelling; microbiome-metabolome interaction networks
    \item[Estimated Completion:] Ongoing longitudinal follow-up; initial
        results published 2025
    \item[Timeline:] Recruitment 2020--2023; initial AI classifier published
        July 2025; longitudinal analysis ongoing
    \item[Publication Medium:] Initial results in high-impact journal (2025);
        longitudinal follow-up expected 2026--2027
    \item[Document Relevance:] Ch.~11 (gut microbiome), Ch.~6 (energy
        metabolism), Ch.~13 (integrative models), Ch.~20 (biomarker research)
\end{description}

BioMapAI is notable for its longitudinal design and AI-first analytical
approach.  Unlike cross-sectional biomarker studies, the 4-year follow-up
can capture disease trajectory and distinguish state markers (fluctuating with
symptoms) from trait markers (stable features of the disease).  The gut
microbiome component strengthens the microbiome--immune--metabolism axis
described in Chapter~11.  The 90\% classification accuracy, while promising,
requires external validation in independent cohorts.

\subsection{ME Association PhD Project (Metabolite and Infection Markers)}
\label{sec:registry-mea-phd}

\begin{description}
    \item[Principal Investigator:] Aleyna Lumsden (PhD candidate); supervised
        across multiple institutions
    \item[Institution:] Rosalind Franklin Institute, Imperial College London,
        and National Phenome Centre, United Kingdom
    \item[Contact/URL:] \url{https://meassociation.org.uk/2025/11/research-new-funding-awarded-to-phd-project-that-will-identify-key-metabolites-and-infection-markers-in-me-cfs/}
    \item[Funder:] ME Association Ramsay Research Fund and UKRI (joint)
    \item[Status:] Active (PhD project started 2025)
    \item[Phase:] Observational, metabolomics and infection biomarkers
    \item[Cohort:] To be determined (likely leveraging existing biobank samples)
    \item[Mechanism/Focus:] Advanced mass spectrometry and informatics to
        identify unknown metabolites and uncover evidence of infection in
        ME/CFS.  Uses facilities across three major UK research centres.
    \item[Primary Outcomes:] Novel metabolite identification; infection
        marker discovery; candidate biomarkers for further validation
    \item[Estimated Completion:] 2028--2029 (typical PhD duration)
    \item[Timeline:] Funding awarded November 2025; project started 2025--2026
    \item[Publication Medium:] PhD thesis and associated journal publications
    \item[Document Relevance:] Ch.~6 (energy metabolism), Ch.~7 (immune
        dysfunction), Ch.~20 (biomarker research)
\end{description}

This PhD project complements the larger Rosetta Stone Study by focusing
specifically on unknown metabolites---compounds not yet characterised in
standard metabolomics panels.  The infection marker component addresses the
long-standing hypothesis that persistent or reactivated infections drive
ME/CFS pathology (Chapter~7).  Access to the Rosalind Franklin Institute's
advanced mass spectrometry facilities gives this project analytical
capabilities beyond typical academic metabolomics studies.

\subsection{MEA--UCL Blood Biomarker Pilot}
\label{sec:registry-mea-ucl}

\begin{description}
    \item[Principal Investigator:] Sophie Hicks (PhD student); supervisors
        Dr.\ Amanda Heslegrave, Dr.\ Michael Zandi, Prof.\ Henrik Zetterberg
    \item[Institution:] UK DRI Fluid Biomarker Laboratory, UCL Institute of
        Neurology, London
    \item[Contact/URL:] ME Association announcement, February 2026
    \item[Registry ID:] Not yet registered (pilot/observational)
    \item[Funder:] ME Association (UK)
    \item[Status:] Active (12-month project, started February 2026)
    \item[Phase:] Observational pilot
    \item[Cohort:] Blood samples from ME/CFS patients (UK ME/CFS Biobank),
        Long Covid patients (UCLH STIMULATE-ICP study), and healthy volunteers
    \item[Mechanism/Focus:] Searches for blood-based biomarkers linking
        specific symptom clusters (brain fog, headaches, muscle pain) to
        measurable biological changes.  Applies neuroimmunology expertise
        from the Zetterberg group (world-leading in fluid biomarkers for
        neurodegeneration, e.g., neurofilament light chain).
    \item[Primary Outcomes:] Candidate blood biomarkers correlated with
        symptom clusters
    \item[Estimated Completion:] $\sim$February 2027
    \item[Timeline:] 12-month pilot; may lead to larger validation study
    \item[Publication Medium:] Not yet announced
    \item[Document Relevance:] Ch.~7 (immune dysfunction),
        Ch.~22 (mechanistic studies)
\end{description}

The Zetterberg laboratory's involvement is notable: their deep expertise in
neurofilament light chain and related neurodegeneration biomarkers has not
previously been applied to ME/CFS.  Cross-disease biomarker candidates
(shared between ME/CFS and neurodegenerative conditions) could emerge from
this work, informing the diagnostic overlap discussion in
Chapter~\ref{sec:cross-disease}.

\subsection{Cornell Cell-Free RNA Liquid Biopsy}
\label{sec:registry-cfrna}

\begin{description}
    \item[Principal Investigator:] Iwijn De Vlaminck (Biomedical Engineering)
        and Maureen Hanson (Molecular Biology and Genetics)
    \item[Institution:] Cornell University, Ithaca, NY
    \item[Contact/URL:] \url{https://doi.org/10.1073/pnas.2507345122}
    \item[Funder:] NIH / institutional
    \item[Status:] Published August 2025
    \item[Phase:] Biomarker discovery (cross-sectional)
    \item[Cohort:] 93 ME/CFS cases + 75 healthy sedentary controls
    \item[Mechanism/Focus:] Profiled circulating cell-free RNA (cfRNA) in
        plasma by RNA sequencing.  Identified over 700 differentially expressed
        transcripts.  Immune cfRNA deconvolution revealed elevated
        plasmacytoid dendritic cell, monocyte, and T~cell-derived cfRNA in
        ME/CFS, with signatures of cytokine signalling and T~cell exhaustion.
    \item[Primary Outcomes:] GLM-LASSO classifier AUC 0.81, accuracy 77\%.
        Not yet diagnostic-grade but substantial advance over prior attempts
    \item[Estimated Completion:] Published
    \item[Publication Medium:] \textit{Proceedings of the National Academy of
        Sciences}
    \item[Document Relevance:] Ch.~7 (immune dysfunction), Ch.~20 (biomarker
        research), Ch.~25 (translational findings)
\end{description}

Cell-free RNA offers a minimally invasive ``liquid biopsy'' approach that
captures tissue-level biology without biopsies.  The 77\% classification
accuracy is insufficient for a standalone diagnostic but demonstrates that
cfRNA carries ME/CFS-relevant biological signal~\cite{Gardella2025cfRNA}.
The T~cell exhaustion and platelet activation signatures align with
independent findings from single-cell transcriptomics
(Section~\ref{sec:registry-tcell-exhaustion}) and proteomics
(Section~\ref{sec:registry-cell-reports-medicine}).  Prospective validation
in a larger, multi-site cohort with disease controls (fibromyalgia,
depression, long COVID) is needed to assess clinical utility.

\subsection{Cell Reports Medicine Multi-Modal Study}
\label{sec:registry-cell-reports-medicine}

\begin{description}
    \item[Principal Investigator:] Multiple collaborating groups
    \item[Institution:] Multi-centre (USA)
    \item[Contact/URL:] \url{https://doi.org/10.1016/j.xcrm.2025.102012}
    \item[Funder:] NIH / institutional
    \item[Status:] Published 2025
    \item[Phase:] Observational, multi-modal
    \item[Cohort:] ME/CFS patients + healthy controls
    \item[Mechanism/Focus:] Integrated energy metabolism assessment, immune
        profiling, and plasma proteomics.  Immune cells from ME/CFS patients
        showed elevated AMP and ADP with reduced ATP/ADP ratio.  NK~cell and
        dendritic cell subset abnormalities.  Vascular dysfunction markers
        including soluble THBS1, VWF, and VASN.
    \item[Primary Outcomes:] Converging evidence across energy, immune, and
        vascular axes; reduced CD56\textsuperscript{low}CD16\textsuperscript{+}
        NK~cells (antiviral), increased plasmacytoid DCs
    \item[Estimated Completion:] Published
    \item[Publication Medium:] \textit{Cell Reports Medicine}
    \item[Document Relevance:] Ch.~6 (energy metabolism), Ch.~7 (immune
        dysfunction), Ch.~10 (cardiovascular), Ch.~13 (integrative models)
\end{description}

This study is noteworthy for its simultaneous assessment of three
pathophysiological axes---energy, immune, and vascular---in the same cohort,
enabling correlation analyses that single-axis studies cannot
perform~\cite{CellReportsMedicine2025MECFS}.  The reduced ATP/ADP ratio in
immune cells provides a mechanistic link between energy metabolism (Ch.~6) and
immune dysfunction (Ch.~7).  The vascular findings (elevated thrombospondin-1,
von Willebrand factor) converge with coagulation abnormalities reported in the
microclot literature~\cite{Nunes2022microclots}.

% =============================================================================
\section{Autoimmunity and Immunomodulation}
\label{sec:registry-autoimmunity}
% =============================================================================

\subsection{Charit\'{e} Immunoadsorption Trial (IA-PACS-CFS)}
\label{sec:registry-charite-ia}

\begin{description}
    \item[Principal Investigator:] Carmen Scheibenbogen, MD, Professor of
        Immunology
    \item[Institution:] Charit\'{e} -- Universit\"{a}tsmedizin Berlin,
        Charit\'{e} Fatigue Centrum, Germany
    \item[Contact/URL:] \url{https://cfc.charite.de/en/clinical_research/nksg/trial_ia_pacs_cfs}
    \item[Registry ID:] NCT05710770
    \item[Funder:] Institutional and ME/CFS Research Foundation
    \item[Status:] Active (double-blind, sham-controlled)
    \item[Phase:] Exploratory interventional trial
    \item[Cohort:] ME/CFS patients (post-COVID and other post-infectious)
        with elevated $\beta_2$-adrenergic receptor autoantibodies
    \item[Mechanism/Focus:] Five immunoadsorption sessions over 10 days
        to remove pathogenic autoantibodies targeting G-protein coupled
        receptors (GPCRs), compared with sham treatment.
        Prior open-label cohort ($n=20$) showed clinical improvement.
    \item[Primary Outcomes:] Clinical improvement at 3 months (Chalder
        Fatigue Scale); autoantibody levels; 6-month follow-up
    \item[Estimated Completion:] Not yet announced
    \item[Timeline:] Open-label cohort 2022--2023 (published January 2025);
        double-blind RCT ongoing
    \item[Publication Medium:] Open-label results in peer-reviewed journal
        (2025); RCT results expected in high-impact immunology journal
    \item[Document Relevance:] Ch.~7 (immune dysfunction---autoantibodies),
        Ch.~10 (autonomic dysfunction), Ch.~19 (integrative approaches)
\end{description}

The Charit\'{e} immunoadsorption programme is the most advanced clinical
investigation of the autoantibody hypothesis in ME/CFS.  Scheibenbogen's group
has systematically built from case reports through an open-label cohort to a
sham-controlled RCT---the methodological progression needed to establish
causality.  The requirement for elevated $\beta_2$-adrenergic receptor
autoantibodies as an inclusion criterion represents a biomarker-stratified
trial design, testing both the treatment and the underlying autoimmune
mechanism simultaneously.  A new 2026-funded neuroimaging substudy will
assess whether thalamocortical hyperconnectivity normalises after
immunoadsorption, linking peripheral autoimmunity to central nervous system
changes.

\subsection{BC007 (Rovunaptabin) Phase II Trials}
\label{sec:registry-bc007}

\begin{description}
    \item[Principal Investigator:] Multiple (industry-sponsored)
    \item[Institution:] Berlin Cures GmbH, Germany (multi-site)
    \item[Contact/URL:] \url{https://clinicaltrials.gov/study/NCT05911009}
    \item[Registry ID:] NCT05911009 (Phase II, $n=114$); earlier Phase IIa
        ($n=29$)
    \item[Funder:] Berlin Cures GmbH (industry)
    \item[Status:] Phase II completed/in analysis
    \item[Phase:] Phase IIa (crossover, $n=29$) and Phase II ($n=114$)
    \item[Cohort:] Post-COVID syndrome patients with functional GPCR
        autoantibodies and predominant fatigue
    \item[Mechanism/Focus:] Rovunaptabin (BC007) is an aptamer that
        neutralises functional autoantibodies against GPCRs.  Unlike
        immunoadsorption (which removes all immunoglobulins non-selectively),
        BC007 targets pathogenic autoantibodies specifically.
    \item[Primary Outcomes:] Fatigue severity; autoantibody neutralisation;
        safety
    \item[Estimated Completion:] Results expected 2026
    \item[Timeline:] Phase IIa enrolled 29 patients; Phase II enrolled 114
        patients; analysis ongoing
    \item[Publication Medium:] Not yet announced; industry-sponsored trials
        typically target high-impact journals
    \item[Document Relevance:] Ch.~7 (immune dysfunction---autoantibodies),
        Ch.~25 (translational findings)
\end{description}

BC007 and immunoadsorption represent complementary approaches to the same
autoantibody hypothesis: immunoadsorption removes antibodies broadly, while
BC007 neutralises specific pathogenic autoantibodies in situ.  The Phase II
trial ($n=114$) is the largest RCT testing autoantibody-targeted therapy in
post-infectious fatigue conditions.  Positive results would validate the
functional GPCR autoantibody hypothesis and open a pharmacological treatment
pathway.  However, the trials enrolled post-COVID patients specifically;
extension to non-COVID ME/CFS cohorts will be needed.

\subsection{ResetME Trial (Daratumumab)}
\label{sec:registry-resetme}

\begin{description}
    \item[Principal Investigator:] \O{}ystein Fl\"{u}ge and Olav Mella
    \item[Institution:] Haukeland University Hospital, Bergen, Norway (with
        Oslo University Hospital)
    \item[Contact/URL:] EudraCT / EU Clinical Trials Register
    \item[Registry ID:] European registry (no NCT number)
    \item[Funder:] ME Fund Norway; institutional
    \item[Status:] Recruiting (opened June 2025)
    \item[Phase:] Phase II, randomised, double-blind, placebo-controlled
        (2:1 ratio)
    \item[Cohort:] 66 participants; moderate-to-severe ME/CFS (CCC criteria),
        minimum 2-year disease duration
    \item[Mechanism/Focus:] Daratumumab is an anti-CD38 monoclonal antibody
        that depletes plasma cells, causing transient reduction of pathogenic
        autoantibodies (including anti-adrenergic and anti-cholinergic receptor
        antibodies).  Follow-on from the open-label KTS9 pilot ($n=10$,
        EudraCT 2022-000281-18), in which 6/10 patients showed marked
        improvement and mean SF-36 Physical Function rose from 25.9 to 55.0.
    \item[Primary Outcomes:] Symptom reduction (SF-36 Physical Function);
        safety
    \item[Estimated Completion:] $\sim$2027 (72-week total participation)
    \item[Timeline:] 26-week treatment period, 72-week follow-up from June
        2025 start
    \item[Publication Medium:] Not yet announced; pilot published in
        \textit{Frontiers in Medicine}, June 2025
    \item[Document Relevance:] Ch.~7 (immune dysfunction---autoantibodies),
        Ch.~10 (adrenergic receptor autoantibodies),
        Ch.~19 (integrative approaches)
\end{description}

ResetME is among the highest-priority ongoing ME/CFS trials.  Fl\"{u}ge and
Mella are the same group that conducted the rituximab trials; daratumumab
targets plasma cells more selectively via CD38.  Unlike rituximab (which
depletes CD20+ B cells), daratumumab directly reduces the antibody-secreting
plasma cells hypothesised to produce pathogenic autoantibodies.  The pilot
response rate (60\%, with near-doubling of physical function scores) is
among the strongest signals in ME/CFS treatment research.  If the RCT
replicates these findings, it would establish the first disease-modifying
treatment for an autoimmune subset of ME/CFS.  Combined with BC007
(Section~\ref{sec:registry-bc007}) and immunoadsorption
(Section~\ref{sec:registry-charite-ia}), this trial represents the third
independent approach to the autoantibody hypothesis.

% =============================================================================
\section{Neuroimaging and Ion Channel Research}
\label{sec:registry-neuro}
% =============================================================================

\subsection{NCNED TRPM3 Ion Channel Programme}
\label{sec:registry-ncned-trpm3}

\begin{description}
    \item[Principal Investigator:] Sonya Marshall-Gradisnik, Professor of
        Immunology
    \item[Institution:] National Centre for Neuroimmunology and Emerging
        Diseases (NCNED), Griffith University, Australia
    \item[Contact/URL:] \url{https://www.griffith.edu.au/research/health/national-centre-neuroimmunology-emerging-diseases}
    \item[Funder:] Institutional and NHMRC (Australia)
    \item[Status:] Active (multiple ongoing studies)
    \item[Phase:] Observational and interventional (LDN trials planned)
    \item[Cohort:] ME/CFS and Long Covid patients (multiple cohorts)
    \item[Mechanism/Focus:] Systematic investigation of TRPM3 ion channel
        dysfunction in ME/CFS.  January 2026 publication confirmed consistent
        TRPM3 malfunction using gold-standard patch-clamp techniques.
        Also conducting 7T MRI studies revealing hippocampal volume
        alterations.  Preparing LDN clinical trials for both ME/CFS
        and Long Covid.
    \item[Primary Outcomes:] TRPM3 channel function; hippocampal volume
        changes; LDN trial endpoints TBD
    \item[Estimated Completion:] Ongoing programme (multiple studies)
    \item[Timeline:] TRPM3 confirmation published January 2026; 7T MRI
        study published February 2025; LDN trials in preparation;
        \$438,000 grant for longitudinal brain imaging (progression study)
    \item[Publication Medium:] Published in \textit{PLOS ONE} and
        specialty journals; future publications expected similarly
    \item[Document Relevance:] Ch.~8 (neurological dysfunction), Ch.~7
        (immune dysfunction---ion channels), Ch.~20 (biomarker research)
\end{description}

The NCNED programme is distinctive for its sustained, systematic focus on a
single molecular target (TRPM3) across multiple methodologies.  The January
2026 patch-clamp confirmation elevates TRPM3 dysfunction from an association
to a consistently reproducible cellular phenotype---a rare achievement in
ME/CFS research.  The planned LDN trials will test whether modulating
neuroinflammation (a proposed upstream driver of ion channel dysfunction)
improves TRPM3 function and clinical outcomes, creating a mechanistic link
between treatment and pathophysiology.  The 7T MRI hippocampal findings add
a neuroanatomical dimension that connects to cognitive complaints reported
by patients.

\subsection{Charit\'{e} Neuroimaging Biomarker Study}
\label{sec:registry-charite-neuroimaging}

\begin{description}
    \item[Principal Investigator:] Carmen Scheibenbogen, MD (PI);
        neuroimaging team at Charit\'{e}
    \item[Institution:] Charit\'{e} -- Universit\"{a}tsmedizin Berlin, Germany
    \item[Contact/URL:] \url{https://mecfs-research.org/en/researchfunding-charite-biomarker2/}
    \item[Funder:] ME/CFS Research Foundation (Germany)
    \item[Status:] Funded; starting January 2026
    \item[Phase:] Observational (neuroimaging biomarker development)
    \item[Cohort:] ME/CFS patients undergoing immunoadsorption treatment
    \item[Mechanism/Focus:] Investigates whether functional brain changes---
        particularly thalamocortical hyperconnectivity---can be modulated by
        immunoadsorption.  Tests whether CNS abnormalities normalise after
        peripheral autoantibody removal, linking the autoimmune and
        neurological hypotheses.
    \item[Primary Outcomes:] Pre/post-treatment fMRI connectivity changes;
        correlation with clinical improvement and autoantibody levels
    \item[Estimated Completion:] Not yet announced (recently funded)
    \item[Timeline:] Funding confirmed late 2025; project starting
        January 2026
    \item[Publication Medium:] Not yet announced
    \item[Document Relevance:] Ch.~8 (neurological dysfunction), Ch.~7
        (immune dysfunction---autoantibodies), Ch.~13 (integrative models)
\end{description}

This study bridges two major ME/CFS hypotheses: autoimmunity and central
nervous system dysfunction.  By imaging the brain before and after
immunoadsorption, it can test whether peripheral autoantibody removal
reverses central connectivity abnormalities---a question that goes beyond
symptom improvement to mechanistic understanding.  If thalamocortical
hyperconnectivity normalises with autoantibody reduction, it would provide
the strongest evidence yet that peripheral immune dysregulation drives
CNS pathology in ME/CFS rather than the reverse.

\subsection{Hyperbaric Oxygen Therapy for ME/CFS}
\label{sec:registry-hbot}

\begin{description}
    \item[Principal Investigator:] Multiple PIs
    \item[Institution:] Multi-site (USA / Germany)
    \item[Contact/URL:] \url{https://doi.org/10.1101/2025.10.29.25339096}
    \item[Registry ID:] NCT06118138 (post-COVID/ME/CFS HBOT study)
    \item[Funder:] ME/CFS Research Foundation / institutional
    \item[Status:] Pilot results published (preprint October 2025);
        larger RCT in planning
    \item[Phase:] Phase 1--2 (open-label pilot)
    \item[Cohort:] $n=30$ ME/CFS patients, 40 HBOT sessions each
    \item[Mechanism/Focus:] Hyperbaric oxygen therapy (2.4~ATA, 100\% O$_2$)
        to address tissue hypoxia, neuroinflammation, and endothelial
        dysfunction.  Pre- and post-treatment fMRI assessed thalamic
        connectivity.
    \item[Primary Outcomes:] SF-36 physical functioning significantly improved.
        Exercise capacity, muscle strength, and information processing speed
        all improved.  Thalamic hyperconnectivity normalised post-HBOT
    \item[Estimated Completion:] Pilot published; RCT timeline TBD
    \item[Publication Medium:] medRxiv preprint; peer review pending
    \item[Document Relevance:] Ch.~8 (neurological dysfunction), Ch.~10
        (cardiovascular), Ch.~19 (integrative approaches)
\end{description}

The normalisation of thalamic hyperconnectivity after HBOT provides the most
direct interventional evidence to date linking brain connectivity abnormalities
to a treatable target in ME/CFS~\cite{HBOT2025mecfs}.  However, the
open-label design and small sample size preclude definitive conclusions;
placebo effects are substantial for fatigue interventions, and a
sham-controlled design (e.g., 1.3~ATA room air) is essential.  The overlap
with brainstem hypoperfusion research strengthens the rationale for
oxygen-targeting approaches.

% =============================================================================
\section{Large-Scale Federally Funded Programmes}
\label{sec:registry-federal}
% =============================================================================

\subsection{RECOVER-TLC Clinical Trials Programme}
\label{sec:registry-recover}

\begin{description}
    \item[Principal Investigator:] Multiple PIs across sites
    \item[Institution:] NIH-funded, multi-site across the United States
    \item[Contact/URL:] \url{https://recovercovid.org/}
    \item[Registry ID:] Multiple NCT numbers (programme-level)
    \item[Funder:] NIH (National Institutes of Health)
    \item[Status:] Active; multiple trials at different stages
    \item[Phase:] Multiple interventional trials (Phase II--III)
    \item[Cohort:] Long Covid patients (with substantial ME/CFS overlap)
    \item[Mechanism/Focus:] The RECOVER-Treating Long COVID (TLC) programme
        includes multiple treatment trials:
        \begin{itemize}
            \item \textbf{REVERSE-LC}: Baricitinib (JAK inhibitor) for
                neurocognitive and cardiopulmonary symptoms; recruiting
            \item \textbf{RECOVER Energize}: Structured pacing for PEM
                ($n=600$); completing December 2025
            \item \textbf{LDN trial}: Low-dose naltrexone for fatigue in
                ages 6--25; enrolment estimated summer 2026
            \item \textbf{GLP-1 agonist trial}: Targeting immune
                dysregulation; enrolment estimated late summer 2026
        \end{itemize}
    \item[Primary Outcomes:] Vary by trial; include symptom severity,
        exercise tolerance, immune markers, quality of life
    \item[Estimated Completion:] Results from first-round trials expected
        by end of 2026
    \item[Timeline:] RECOVER programme initiated 2021; TLC trials began
        2024--2025; second-round trials planning 2026
    \item[Publication Medium:] High-impact journals; all results expected
        to be open access (NIH-funded)
    \item[Document Relevance:] Ch.~7 (immune dysfunction), Ch.~10
        (cardiovascular), Ch.~17 (pacing), Ch.~14
        \S\ref{sec:cross-disease} (cross-disease)
\end{description}

RECOVER-TLC is the largest federally funded treatment trial programme
relevant to ME/CFS.  While focused on Long Covid, the substantial clinical
overlap means positive results are directly translatable.  The baricitinib
(JAK inhibitor) trial is particularly noteworthy given the JAK/STAT pathway
abnormalities identified in ME/CFS (Chapter~7) and the EpiSwitch epigenetic
findings.  The structured pacing trial (RECOVER Energize, $n=600$) will
provide the first large-scale evidence for pacing---the most widely
recommended non-pharmacological intervention in ME/CFS---in a rigorous
clinical trial framework.  The paediatric LDN trial fills a critical gap,
as most ME/CFS treatment evidence excludes children and adolescents.

% =============================================================================
\section{Clinical Trials (Interventional)}
\label{sec:registry-interventional}
% =============================================================================

\subsection{LIFT Trial (Low-Dose Naltrexone and Pyridostigmine)}
\label{sec:registry-lift}

\begin{description}
    \item[Principal Investigator:] David Systrom, MD
    \item[Institution:] Brigham and Women's Hospital, Harvard Medical School,
        United States
    \item[Contact/URL:] \url{https://solvecfs.org/exploring-current-me-cfs-and-long-covid-clinical-trials-on-national-clinical-trials-day/}
    \item[Funder:] Solve ME/CFS Initiative (Ramsay Research Grant)
    \item[Status:] Recruiting
    \item[Phase:] Clinical trial (interventional)
    \item[Cohort:] ME/CFS patients (specific $n$ to be confirmed)
    \item[Mechanism/Focus:] Tests whether low-dose naltrexone (LDN) and
        pyridostigmine reduce symptoms and improve quality of life.  LDN
        targets neuroinflammation via glial cell modulation; pyridostigmine
        is an acetylcholinesterase inhibitor used for autonomic dysfunction.
    \item[Primary Outcomes:] Symptom severity; quality of life measures
    \item[Estimated Completion:] Not yet announced
    \item[Timeline:] Recruitment ongoing as of 2026
    \item[Publication Medium:] Not yet announced
    \item[Document Relevance:] Ch.~16 (supplements and nutraceuticals),
        Ch.~10 (cardiovascular/autonomic), Ch.~17 (lifestyle interventions)
\end{description}

The LIFT trial addresses two of the most widely used off-label treatments in
ME/CFS clinical practice.  LDN has accumulated substantial anecdotal and
preliminary evidence (Chapter~16) but lacks rigorous RCT data in ME/CFS
specifically.  Pyridostigmine targets the preload failure mechanism identified
in invasive cardiopulmonary exercise testing by Systrom's own group
(Chapter~10).  A positive result could provide the first high-quality evidence
base for two treatments already in widespread empirical use.

\subsection{UAB LDN Dose-Finding Study}
\label{sec:registry-uab-ldn}

\begin{description}
    \item[Principal Investigator:] Jarred Younger, PhD
    \item[Institution:] University of Alabama at Birmingham, United States
    \item[Contact/URL:] \url{https://clinicaltrials.gov/study/NCT07285473}
    \item[Registry ID:] NCT07285473
    \item[Funder:] Institutional
    \item[Status:] Recruiting (remote/decentralised, US-wide)
    \item[Phase:] Dose-finding study
    \item[Cohort:] Adults with ME/CFS (specific $n$ TBD)
    \item[Mechanism/Focus:] Dose-finding study for low-dose naltrexone in
        ME/CFS.  Distinct from the LIFT trial (which combines LDN with
        pyridostigmine, Section~\ref{sec:registry-lift}) and the RECOVER
        paediatric LDN trial (Section~\ref{sec:registry-recover}): this
        study focuses on determining optimal adult LDN dosing specifically
        for ME/CFS.  Remote design enables broad US enrolment.
    \item[Primary Outcomes:] Dose-response for symptom reduction;
        safety and tolerability
    \item[Estimated Completion:] Not yet announced
    \item[Timeline:] Recruitment ongoing as of 2026
    \item[Publication Medium:] Not yet announced
    \item[Document Relevance:] Ch.~16 (supplements---LDN),
        Ch.~7 (immune dysfunction---microglial hypothesis)
\end{description}

Younger's laboratory has prior immunology work on LDN in ME/CFS
(NCT02965768).  While the LIFT trial tests LDN in combination,
this study isolates LDN's independent effect and seeks to establish
the optimal dose---a fundamental gap in the current evidence base.
Most clinicians prescribe 1.5--4.5~mg empirically; systematic dose-finding
data could substantially improve clinical practice.

\subsection{Low-Dose Rapamycin in ME/CFS and Long Covid}
\label{sec:registry-rapamycin}

\begin{description}
    \item[Principal Investigator:] Multi-site collaboration
    \item[Institution:] Simmaron Research Institute, Bateman Horne Center,
        CDC, and Mayo Clinic, United States
    \item[Contact/URL:] \url{https://batemanhornecenter.org/promising-clinical-trials/}
    \item[Funder:] Multi-institutional
    \item[Status:] Active
    \item[Phase:] Clinical trial (interventional)
    \item[Cohort:] ME/CFS and Long Covid patients (specific $n$ TBD)
    \item[Mechanism/Focus:] Rapamycin (sirolimus) is an mTOR inhibitor with
        immunomodulatory and anti-inflammatory properties.  Targets mTOR
        pathway dysregulation implicated in ME/CFS immune and metabolic
        abnormalities.
    \item[Primary Outcomes:] Symptom improvement; immune marker changes;
        safety and tolerability
    \item[Estimated Completion:] Not yet announced
    \item[Timeline:] Study active as of 2025--2026
    \item[Publication Medium:] Not yet announced
    \item[Document Relevance:] Ch.~7 (immune dysfunction), Ch.~6 (energy
        metabolism), Ch.~19 (integrative approaches)
\end{description}

The rapamycin trial is notable for its multi-institutional design (Simmaron,
Bateman Horne, CDC, Mayo Clinic), which lends credibility and ensures
phenotyping rigour.  mTOR pathway dysregulation has been reported in ME/CFS
(Chapter~7) and connects immune activation to metabolic dysfunction
(Chapter~6).  This trial directly tests a mechanistic hypothesis rather than
treating symptoms empirically, making it particularly informative regardless of
outcome.

\subsection{Lumbrokinase for ME/CFS and Long Covid}
\label{sec:registry-lumbrokinase}

\begin{description}
    \item[Principal Investigator:] David Putrino, PhD
    \item[Institution:] Abilities Research Center, Mount Sinai Icahn School
        of Medicine, United States
    \item[Contact/URL:] \url{https://solvecfs.org/exploring-current-me-cfs-and-long-covid-clinical-trials-on-national-clinical-trials-day/}
    \item[Funder:] Institutional
    \item[Status:] Recruiting
    \item[Phase:] Clinical trial (interventional)
    \item[Cohort:] Adults with Long Covid, post-treatment Lyme disease
        syndrome, and ME/CFS
    \item[Mechanism/Focus:] Lumbrokinase is a fibrinolytic enzyme targeting
        microclot pathology.  The microclot hypothesis proposes that
        persistent fibrin amyloid microclots impair microcirculation,
        contributing to tissue hypoxia and fatigue.
    \item[Primary Outcomes:] Fibrinolytic markers; symptom scores; exercise
        tolerance
    \item[Estimated Completion:] Not yet announced
    \item[Timeline:] Recruitment ongoing as of 2026
    \item[Publication Medium:] Not yet announced
    \item[Document Relevance:] Ch.~10 (cardiovascular dysfunction),
        Ch.~19 (integrative approaches)
\end{description}

The microclot hypothesis has generated significant interest since 2021 but
lacks interventional trial data.  This study directly tests whether
fibrinolytic treatment improves outcomes, providing a causal test of the
hypothesis.  The inclusion of three post-infectious conditions (Long Covid,
post-Lyme, ME/CFS) mirrors the Rosetta Stone Study's cross-disease approach,
potentially revealing whether microclot pathology is a shared mechanism.

\subsection{Hydrogen Water for ME/CFS}
\label{sec:registry-hydrogen-water}

\begin{description}
    \item[Principal Investigator:] Fred Friedberg, PhD
    \item[Institution:] Stony Brook University School of Medicine,
        United States
    \item[Contact/URL:] \url{https://solvecfs.org/exploring-current-me-cfs-and-long-covid-clinical-trials-on-national-clinical-trials-day/}
    \item[Funder:] Solve ME/CFS Initiative
    \item[Status:] Recruiting
    \item[Phase:] Clinical trial, dosing study
    \item[Cohort:] ME/CFS patients (specific $n$ TBD)
    \item[Mechanism/Focus:] Hydrogen-rich water as an antioxidant
        intervention.  Molecular hydrogen ($\mathrm{H}_2$) selectively
        scavenges hydroxyl radicals and may reduce oxidative stress and
        neuroinflammation.
    \item[Primary Outcomes:] Fatigue severity; stress markers; physical
        function
    \item[Estimated Completion:] Not yet announced
    \item[Timeline:] Recruitment ongoing as of 2026
    \item[Publication Medium:] Not yet announced
    \item[Document Relevance:] Ch.~6 (energy metabolism---oxidative stress),
        Ch.~16 (supplements and nutraceuticals)
\end{description}

This dosing study addresses a practical gap: even if hydrogen water has
antioxidant benefits, optimal dosing for ME/CFS has not been established.
The oxidative stress hypothesis in ME/CFS (Chapter~6) provides a mechanistic
rationale, though the evidence base for molecular hydrogen in any chronic
condition remains preliminary.  Results will be most informative if they
include objective oxidative stress biomarkers alongside subjective symptom
measures.

\subsection{RESTORE-ME: Oxaloacetate for ME/CFS Fatigue}
\label{sec:registry-restore-me}

\begin{description}
    \item[Principal Investigator:] Alan Cash, PhD
    \item[Institution:] Terra Biological LLC, United States (multi-site RCT)
    \item[Contact/URL:] \url{https://clinicaltrials.gov/study/NCT05273372}
    \item[Registry ID:] NCT05273372
    \item[Funder:] Terra Biological LLC (industry)
    \item[Status:] Results published (November 2024)
    \item[Phase:] Phase II, randomised, double-blind, placebo-controlled
    \item[Cohort:] 82 ME/CFS patients (2,000\,mg oxaloacetate vs control
        daily for 3 months)
    \item[Mechanism/Focus:] Oxaloacetate (OAA) is a mitochondrial metabolite
        that supports the citric acid cycle and NAD$^+$/NADH ratio.
        Targets metabolic dysfunction by supplementing a key TCA cycle
        intermediate.
    \item[Primary Outcomes:] Fatigue reduced $>$25\% from baseline in
        OAA group ($p=0.0039$ vs control); cognitive improvement at
        day~60 ($p=0.034$); modest increase in uptime at day~30
    \item[Estimated Completion:] Completed
    \item[Timeline:] RCT completed 2024; results published November 2024
        in \textit{Frontiers in Neurology}; cognitive follow-up published
        2025
    \item[Publication Medium:] \textit{Frontiers in Neurology} (open access)
    \item[Document Relevance:] Ch.~6 (energy metabolism---TCA cycle),
        Ch.~16 (supplements and nutraceuticals)
\end{description}

RESTORE-ME is one of very few completed RCTs showing statistically significant
fatigue reduction in ME/CFS.  The mechanistic rationale is direct: oxaloacetate
is an obligate TCA cycle intermediate, and ME/CFS metabolomic studies have
consistently identified TCA cycle disruption (Chapter~6).  The 25\% fatigue
reduction is clinically meaningful, though the modest sample size ($n=82$)
and industry sponsorship warrant independent replication.  The cognitive
improvement findings strengthen the case that metabolic support can address
both physical and cognitive symptoms.

\subsection{Stellate Ganglion Block for ME/CFS}
\label{sec:registry-sgb}

\begin{description}
    \item[Principal Investigator:] Deborah Duricka, PhD; Lipeng Liu, MD
    \item[Institution:] Multi-site, United States
    \item[Contact/URL:] \url{https://doi.org/10.1080/21641846.2025.2455876}
    \item[Funder:] Solve ME/CFS Initiative (Ramsay Research Grant)
    \item[Status:] Pilot study published (February 2025); larger trials needed
    \item[Phase:] Prospective cohort pilot ($n=10$)
    \item[Cohort:] 10 patients meeting both WHO Long Covid and IOM ME/CFS
        criteria
    \item[Mechanism/Focus:] Sequential bilateral stellate ganglion blocks
        (local anaesthetic injection near the cervical sympathetic chain)
        administered over 3 consecutive weeks.  Hypothesised to ``reset''
        sympathetic/parasympathetic balance of the autonomic nervous system,
        increasing regional blood flow and reducing sympathetic hyperactivation.
    \item[Primary Outcomes:] Significant improvements in cognition, physical
        and social functioning, vitality; reduced orthostatic intolerance
        and POTS symptoms.  Cortisol levels did not change.
    \item[Estimated Completion:] Pilot completed; larger RCT recommended
    \item[Timeline:] Pilot published February 2025; larger trial planning
        stage
    \item[Publication Medium:] \textit{Fatigue: Biomedicine, Health \&
        Behavior} (2025)
    \item[Document Relevance:] Ch.~10 (cardiovascular/autonomic dysfunction),
        Ch.~8 (neurological dysfunction), Ch.~19 (integrative approaches)
\end{description}

The stellate ganglion block is a well-established procedure in pain medicine
repurposed for autonomic dysfunction.  The pilot results are notable for
improvement across both autonomic (OI, POTS) and cognitive domains,
suggesting a shared upstream mechanism.  The ``autonomic reset'' hypothesis
aligns with evidence of sympathetic hyperactivation in ME/CFS (Chapter~10).
The very small sample ($n=10$) limits conclusions, but the effect sizes and
the procedure's established safety profile justify a larger sham-controlled
RCT.

\subsection{Transcutaneous Vagus Nerve Stimulation (tVNS) for Long Covid/ME/CFS}
\label{sec:registry-tvns}

\begin{description}
    \item[Principal Investigator:] Multiple investigators (several parallel
        studies)
    \item[Institution:] Multi-site (US and European studies)
    \item[Contact/URL:] \url{https://clinicaltrials.gov/study/NCT06585254}
    \item[Registry ID:] NCT06585254 (US dosing study); additional European
        trials ongoing
    \item[Funder:] Multiple (Solve ME/CFS Initiative, institutional)
    \item[Status:] Recruiting (multiple trials)
    \item[Phase:] Interventional (parameter optimisation and efficacy)
    \item[Cohort:] Long Covid and ME/CFS patients fulfilling CFS criteria;
        prior studies $n=10$--25 per trial
    \item[Mechanism/Focus:] Non-invasive auricular stimulation of the vagus
        nerve to activate the cholinergic anti-inflammatory pathway (CAP).
        Targets neuroinflammation, autonomic imbalance, and systemic
        inflammation.  Prior pilot studies showed improvements in cognition,
        anxiety, depression, sleep, and autonomic scores (COMPASS-31).
    \item[Primary Outcomes:] Optimal stimulation parameters; health-related
        quality of life; fatigue severity; autonomic function
    \item[Estimated Completion:] Varies by trial; results expected 2026--2027
    \item[Timeline:] Multiple pilot studies published 2024--2025; current
        trials optimising parameters and assessing efficacy in larger cohorts
    \item[Publication Medium:] \textit{Frontiers in Neurology}, \textit{PLOS
        ONE}, and others (prior studies published open access)
    \item[Document Relevance:] Ch.~8 (neurological dysfunction), Ch.~10
        (autonomic dysfunction), Ch.~7 (immune dysfunction---inflammation)
\end{description}

Vagus nerve stimulation represents a compelling intersection of neurology and
immunology: the cholinergic anti-inflammatory pathway provides a mechanism
through which neural stimulation can modulate systemic inflammation.  The
consistency of improvements across multiple small pilot studies (cognition,
autonomic function, sleep) is encouraging, though no large RCT has yet been
completed.  The non-invasive nature, low side-effect profile, and consumer
device availability (e.g., Nurosym) make tVNS a practical intervention if
efficacy is confirmed.  The current parameter-optimisation trials are a
necessary step before definitive efficacy studies.

\subsection{Masitinib for Severe Mast Cell Activation Syndrome}
\label{sec:registry-masitinib}

\begin{description}
    \item[Principal Investigator:] Multiple investigators
    \item[Institution:] AB Science (sponsor), multi-site
    \item[Contact/URL:] \url{https://clinicaltrials.gov/study/NCT05449444}
    \item[Registry ID:] NCT05449444
    \item[Funder:] AB Science (industry)
    \item[Status:] Phase 2 (active)
    \item[Phase:] Phase 2 clinical trial
    \item[Cohort:] Patients with severe mast cell activation syndrome (MCAS)
    \item[Mechanism/Focus:] Masitinib is a tyrosine kinase inhibitor that
        targets mast cell activity directly, unlike antihistamines that only
        address downstream mediator release.  MCAS is identified in an
        estimated 50\% of ME/CFS patients by some specialists.
    \item[Primary Outcomes:] MCAS symptom severity; safety and tolerability
    \item[Estimated Completion:] Not yet announced
    \item[Timeline:] Trial active as of 2025--2026
    \item[Publication Medium:] Not yet announced
    \item[Document Relevance:] Ch.~7 (immune dysfunction---mast cells),
        Ch.~19 (integrative approaches)
\end{description}

Mast cell activation has emerged as a significant comorbidity in ME/CFS
subsets, with some clinicians reporting MCAS in up to half of their ME/CFS
patients.  A 2026 cromolyn sodium case series from Dutch and Johns Hopkins
investigators showed improvement in ME/CFS-related symptoms alongside MCAS
management.  Masitinib's upstream mechanism (direct tyrosine kinase
inhibition rather than antihistamine blockade) could provide more complete
mast cell stabilisation.  If effective, this trial could open a treatment
pathway for the MCAS subset of ME/CFS patients who respond poorly to
conventional antihistamine regimens.

\subsection{Cornell NAC--Glutathione Brain Study}
\label{sec:registry-cornell-nac}

\begin{description}
    \item[Principal Investigator:] Dikoma Shungu, PhD
    \item[Institution:] Weill Cornell Medical College, United States
    \item[Contact/URL:] \url{https://clinicaltrials.gov/study/NCT04542161}
    \item[Registry ID:] NCT04542161
    \item[Funder:] NINDS (National Institute of Neurological Disorders and
        Stroke)
    \item[Status:] Recruiting (Phase 2)
    \item[Phase:] Phase 2, double-blind, placebo-controlled, randomised
    \item[Cohort:] 60 participants (20 per arm: 0~mg / 900~mg / 3600~mg
        daily NAC), ages 21--60, baseline brain GSH at or below predefined
        cutoff, primary ME/CFS diagnosis
    \item[Mechanism/Focus:] Tests N-acetylcysteine (NAC) as a glutathione
        (GSH) precursor to restore cortical GSH reserves depleted by
        oxidative stress.  Uses magnetic resonance spectroscopy (MRS) to
        measure brain GSH levels as a direct target engagement biomarker.
        Three-arm dose-response design.
    \item[Primary Outcomes:] Brain GSH levels (MRS); plasma oxidative stress
        markers
    \item[Estimated Completion:] $\sim$2026
    \item[Timeline:] Recruitment ongoing
    \item[Publication Medium:] Not yet announced
    \item[Document Relevance:] Ch.~6 (energy metabolism---oxidative stress),
        Ch.~22 (mechanistic studies)
\end{description}

This trial is notable for using MRS as both a mechanistic measurement and
treatment-response biomarker in a single tool.  If NAC restores brain GSH
levels and correlates with symptom improvement, it would establish oxidative
stress as a druggable target in ME/CFS and validate a non-invasive brain
biomarker for treatment monitoring.  The three-arm dose-response design
addresses the fundamental question of whether standard supplement doses
(900~mg) achieve meaningful CNS penetration.

\subsection{NSU Directed Probiotics for ME/CFS}
\label{sec:registry-nsu-probiotics}

\begin{description}
    \item[Principal Investigator:] Nancy Klimas, MD
    \item[Institution:] Nova Southeastern University, Institute for
        Neuro-Immune Medicine, United States
    \item[Contact/URL:] \url{https://clinicaltrials.gov/study/NCT06211062}
    \item[Registry ID:] NCT06211062
    \item[Funder:] Institutional
    \item[Status:] Recruiting (Phase 2)
    \item[Phase:] Phase 2, randomised, placebo-controlled
    \item[Cohort:] 100 participants (25 per arm $\times$ 4 arms: ME/CFS
        with/without IBS $\times$ active/placebo), ages 45--70, Canadian
        Consensus Criteria
    \item[Mechanism/Focus:] Tests i3.1 probiotic (Floradapt Intensive GI) to
        reduce GI inflammation, normalise gut--brain axis, and reset
        microbiome.  Four-arm design allows interaction analysis between
        IBS comorbidity and treatment response.
    \item[Primary Outcomes:] GI inflammation markers; IBS severity score;
        ME/CFS global function; assessed at 8 and 12 weeks
    \item[Estimated Completion:] $\sim$2026
    \item[Timeline:] Recruitment ongoing
    \item[Publication Medium:] Not yet announced
    \item[Document Relevance:] Ch.~11 (gut microbiome---primary),
        Ch.~7 (immune dysfunction---gut--immune axis)
\end{description}

This is the first ME/CFS-specific probiotic RCT with IBS stratification.
The Klimas group's neuro-immune medicine expertise and the four-arm design
(separating ME/CFS-only from ME/CFS+IBS) could determine whether gut-targeted
interventions benefit all ME/CFS patients or specifically those with GI
comorbidity---a distinction with direct clinical implications for the
microbiome hypothesis discussed in Chapter~11.

\subsection{ADDRESS-LC (Bezisterim for Long Covid Neurocognition)}
\label{sec:registry-address-lc}

\begin{description}
    \item[Principal Investigator:] Multi-site (BioVie Inc.)
    \item[Institution:] Multi-centre, United States (DoD-funded)
    \item[Contact/URL:] \url{https://clinicaltrials.gov/study/NCT06847191}
    \item[Registry ID:] NCT06847191
    \item[Funder:] Department of Defense (US)
    \item[Status:] Recruiting (first patient enrolled May 2025)
    \item[Phase:] Phase 2, multi-centre, double-blind, RCT
    \item[Cohort:] 200 adults with Long Covid, enriched for cognitive
        impairment and fatigue, symptoms $\geq$3 months
    \item[Mechanism/Focus:] Bezisterim (NE3107) is an orally bioavailable,
        blood--brain barrier-permeable insulin sensitiser and selective
        NF-$\kappa$B/ERK modulator.  Blocks TNF-$\alpha$ and
        neuroinflammation without global immunosuppression.  Targets the
        neuroinflammation hypothesis for brain fog.  Same drug being
        explored in Alzheimer's and Parkinson's.
    \item[Primary Outcomes:] Cogstate Cognition battery (cognitive
        performance); secondary: fatigue
    \item[Estimated Completion:] Data expected H1 2026
    \item[Timeline:] 84-day treatment period
    \item[Publication Medium:] Not yet announced
    \item[Document Relevance:] Ch.~7 (immune dysfunction---neuroinflammation),
        Ch.~22 (mechanistic studies---Long Covid overlap)
\end{description}

ADDRESS-LC is notable for targeting neuroinflammation via a mechanism
distinct from LDN: selective NF-$\kappa$B modulation rather than glial
cell modulation.  The blood--brain barrier permeability and dual
insulin-sensitising/anti-inflammatory action address two pathophysiological
axes simultaneously.  Cross-disease application (Alzheimer's, Parkinson's,
Long Covid) could accelerate regulatory pathways and provide comparative
data on shared neuroinflammatory mechanisms.

\subsection{LC-REVITALIZE (Upadacitinib + Pirfenidone)}
\label{sec:registry-lc-revitalize}

\begin{description}
    \item[Principal Investigator:] Douglas Fraser, MD PhD
    \item[Institution:] Western University, London, Ontario, Canada;
        Schmidt Initiative for Long COVID (SILC); multi-national
    \item[Contact/URL:] \url{https://clinicaltrials.gov/study/NCT06928272}
    \item[Registry ID:] NCT06928272
    \item[Funder:] Schmidt Initiative for Long COVID (SILC)
    \item[Status:] Recruiting (Phase 3)
    \item[Phase:] Phase 3, double-blind, placebo-controlled, multi-arm
        platform
    \item[Cohort:] $\sim$348 adults with Long Covid (prior confirmed
        SARS-CoV-2, symptoms $\geq$3 months); sites in Brazil, Canada,
        Italy, Uganda, USA, Zambia
    \item[Mechanism/Focus:] Repurposed drug platform trial.  Upadacitinib
        (JAK1 inhibitor, approved for rheumatoid arthritis) and pirfenidone
        (anti-fibrotic, approved for lung disease) identified via AI
        proteomics screen of 5,400 blood proteins from 1,028 participants.
    \item[Primary Outcomes:] Five symptom domains: fatigue, breathlessness,
        cognition, musculoskeletal pain, circulation
    \item[Estimated Completion:] $\sim$2027
    \item[Timeline:] Recruitment ongoing across 6 countries
    \item[Publication Medium:] Not yet announced
    \item[Document Relevance:] Ch.~7 (immune dysfunction---JAK pathway),
        Ch.~13 (integrative models),
        Ch.~22 (mechanistic studies---Long Covid overlap)
\end{description}

LC-REVITALIZE is the first Long Covid treatment trial to use AI-driven
proteomics for drug selection, screening 5,400 proteins to identify
candidate interventions.  The JAK1 inhibitor arm complements the RECOVER-TLC
baricitinib trial (Section~\ref{sec:registry-recover}), testing a more
selective JAK inhibitor.  The multi-national platform design across six
countries enables geographic and genetic diversity rarely achieved in
ME/CFS-adjacent trials.  If positive, the AI proteomics drug-selection
methodology could be applied directly to ME/CFS cohorts.

\subsection{Viral Persistence Trials (Paxlovid/STOP-PASC)}
\label{sec:registry-viral-persistence}

\begin{description}
    \item[Principal Investigator:] Multiple (NIH-funded)
    \item[Institution:] Multi-site across the United States
    \item[Contact/URL:] \url{https://clinicaltrials.gov/study/NCT05576662}
    \item[Registry ID:] NCT05576662 (STOP-PASC); PAX~LC (separate trial)
    \item[Funder:] NIH
    \item[Status:] Completed; results published 2025
    \item[Phase:] Phase 2, randomised, double-blind, placebo-controlled
    \item[Cohort:] Adults with established Long Covid (STOP-PASC and PAX~LC
        trials)
    \item[Mechanism/Focus:] Nirmatrelvir/ritonavir (Paxlovid) for 15 days
        to test whether antiviral treatment targeting SARS-CoV-2 viral
        persistence improves Long Covid symptoms.  Both STOP-PASC and
        PAX~LC trials showed \textbf{no therapeutic benefit}, despite
        the viral persistence hypothesis.
    \item[Primary Outcomes:] No significant improvement in Long Covid
        symptoms vs placebo in either trial
    \item[Estimated Completion:] Completed
    \item[Timeline:] Results from STOP-PASC (systems immunology analysis)
        published December 2025; PAX~LC results published in
        \textit{The Lancet Infectious Diseases} (2025)
    \item[Publication Medium:] \textit{The Lancet Infectious Diseases};
        \textit{medRxiv} (systems immunology analysis)
    \item[Document Relevance:] Ch.~7 (immune dysfunction), Ch.~14
        \S\ref{sec:cross-disease} (cross-disease), Ch.~24 (controversies)
\end{description}

The negative results from STOP-PASC and PAX~LC are as informative as positive
findings would have been.  If SARS-CoV-2 viral persistence were the primary
driver of Long Covid symptoms, a 15-day course of a potent antiviral should
have produced measurable improvement.  The failure to do so suggests either
that viral persistence is not the dominant mechanism in established Long Covid,
that the treatment duration or timing was insufficient, or that downstream
consequences of initial infection (immune dysregulation, autoimmunity,
microbiome disruption) persist independently of ongoing viral replication.
This has direct implications for ME/CFS, where post-infectious mechanisms
similarly outlast the triggering infection.

\subsection{CHIIME (Chronic Infections and Inflammation in ME/CFS)}
\label{sec:registry-chiime}

\begin{description}
    \item[Principal Investigator:] UCSF / PolyBio Research Foundation network
        (LIINC collaborative); Dr.\ John Chia (Enterovirus Medical Research
        Center) as key partner
    \item[Institution:] University of California San Francisco, United States
    \item[Contact/URL:] \url{https://clinicaltrials.gov/study/NCT07227441}
    \item[Registry ID:] NCT07227441
    \item[Funder:] PolyBio Research Foundation
    \item[Status:] Recruiting (start date January 2026)
    \item[Phase:] Observational, prospective, longitudinal
    \item[Cohort:] ME/CFS patients with pre-2019 onset only (excludes Long
        Covid onset to avoid confounding)
    \item[Mechanism/Focus:] Investigates viral persistence (particularly
        enterovirus tissue persistence in gut, muscle, brain),
        T-cell neuroinflammation, and longitudinal biological signatures.
        Gut tissue collected from a subset via biopsy to search for
        enterovirus RNA/protein.  Uses whole-body PET imaging and
        single-nuclei RNA sequencing.
    \item[Primary Outcomes:] Longitudinal biomarker trajectories; enterovirus
        persistence confirmation; neuroimmune phenotyping
    \item[Estimated Completion:] September 2030
    \item[Timeline:] Multi-year longitudinal follow-up
    \item[Publication Medium:] Not yet announced
    \item[Document Relevance:] Ch.~7 (immune dysfunction---viral persistence),
        Ch.~11 (gut microbiome), Ch.~13 (integrative models),
        Ch.~22 (mechanistic studies)
\end{description}

CHIIME represents a paradigm shift from serology-based viral persistence
research to direct tissue-level detection.  By collecting gut, muscle, and
CNS-adjacent tissue, this study can directly test whether enterovirus RNA
and protein persist in tissues of ME/CFS patients---moving beyond the
indirect evidence that has characterised viral persistence research for
decades.  The pre-2019 onset criterion ensures that findings apply to
classical ME/CFS rather than post-COVID presentations, addressing a
critical gap in the field.

% =============================================================================
\section{Virology and Post-Infectious Mechanisms}
\label{sec:registry-virology}
% =============================================================================

\subsection{Herpesvirus Reactivation and Anti-dUTPase Antibodies}
\label{sec:registry-herpesvirus-reactivation}

\begin{description}
    \item[Principal Investigator:] Mar{\'\i}a Palomo and collaborators
    \item[Institution:] Multi-centre (Spain / USA)
    \item[Contact/URL:] \url{https://doi.org/10.1002/jmv.70769}
    \item[Funder:] Institutional
    \item[Status:] Published 2026
    \item[Phase:] Observational (case-control)
    \item[Cohort:] ME/CFS patients + healthy controls
    \item[Mechanism/Focus:] Measured antibodies to EBV, HHV-6, and VZV, plus
        anti-dUTPase IgG (a marker of active herpesvirus replication).
        72.5\% of ME/CFS patients simultaneously co-expressed antibodies to
        multiple herpesviruses versus 31\% of controls.  Heightened
        herpesvirus antibody levels clustered with moderate-to-severe fatigue.
    \item[Primary Outcomes:] Multi-herpesvirus co-reactivation as hallmark of
        post-infectious ME/CFS; anti-dUTPase IgG as fatigue-associated marker
    \item[Estimated Completion:] Published
    \item[Publication Medium:] \textit{Journal of Medical Virology}
    \item[Document Relevance:] Ch.~7 (immune dysfunction), Ch.~9
        (infection and triggers), Ch.~20 (biomarker research)
\end{description}

The simultaneous co-reactivation of three herpesviruses in over 70\% of
patients suggests systemic immune surveillance failure rather than
reactivation of a single pathogen~\cite{Palomo2026herpesvirus}.  The
anti-dUTPase antibody is mechanistically interesting because dUTPase is an
enzyme conserved across herpesviruses that can directly activate innate
immune signalling.  If validated as a stratification marker, it could
identify a ``viral reactivation'' subtype amenable to antiviral or
immunomodulatory therapy.

\subsection{Post-IM Longitudinal Cohort: Seven-Year Follow-Up}
\label{sec:registry-post-im-cohort}

\begin{description}
    \item[Principal Investigator:] Ben Z.\ Katz, MD
    \item[Institution:] Northwestern University / multi-site (USA)
    \item[Contact/URL:] \url{https://doi.org/10.3389/fmed.2026.1676628}
    \item[Funder:] NIH
    \item[Status:] Published 2026
    \item[Phase:] Observational, prospective longitudinal cohort
    \item[Cohort:] 4,501 college students enrolled before developing
        infectious mononucleosis (IM); followed for 7~years
    \item[Mechanism/Focus:] Uniquely prospective design: baseline data
        collected at least 6~weeks before IM onset.  Tracks natural history
        from pre-infection through ME/CFS development, enabling
        identification of pre-existing risk factors.
    \item[Primary Outcomes:] ME/CFS incidence: 13\% at 6~months, declining to
        $\sim$4\% at 24~months post-IM.  Risk factor identification at 7-year
        endpoint
    \item[Estimated Completion:] Published
    \item[Publication Medium:] \textit{Frontiers in Medicine}
    \item[Document Relevance:] Ch.~9 (infection and triggers), Ch.~23
        (epidemiology), Ch.~14 \S\ref{sec:cross-disease} (post-viral overlap)
\end{description}

This is one of the few truly prospective studies in ME/CFS, with pre-infection
baseline data that enables causal inference about risk factors---a design
strength shared with the DecodeME genetic study
(Section~\ref{sec:registry-decodeme}).  The declining incidence from 13\% to
4\% over two years suggests spontaneous recovery in most post-IM ME/CFS
cases, but the persistent 4\% represents a ``locked-in'' subset of
particular clinical and research interest~\cite{Katz2026mononucleosis7yr}.

\subsection{RECOVER ME/CFS Incidence Study}
\label{sec:registry-recover-incidence}

\begin{description}
    \item[Principal Investigator:] RECOVER Consortium
    \item[Institution:] NIH multi-site (USA)
    \item[Contact/URL:] \url{https://doi.org/10.1007/s11606-024-09290-9}
    \item[Funder:] NIH RECOVER Initiative
    \item[Status:] Published 2025
    \item[Phase:] Observational (adult cohort)
    \item[Cohort:] Large-scale RECOVER-Adult observational study
    \item[Mechanism/Focus:] Determined incidence and prevalence of ME/CFS
        following SARS-CoV-2 infection.  Applies established diagnostic
        criteria to the largest prospective post-COVID cohort.
    \item[Primary Outcomes:] Confirmed that ME/CFS incidence increases
        following SARS-CoV-2 infection; provides population-level estimates
    \item[Estimated Completion:] Published
    \item[Publication Medium:] \textit{Journal of General Internal Medicine}
    \item[Document Relevance:] Ch.~9 (infection and triggers), Ch.~23
        (epidemiology), Ch.~14 \S\ref{sec:cross-disease}
\end{description}

This study provides the strongest epidemiological evidence to date that
SARS-CoV-2 infection increases ME/CFS incidence, confirming a pattern
previously observed with EBV, Ross River virus, and \textit{Coxiella
burnetii}~\cite{RECOVER2025}.  The RECOVER infrastructure's
scale enables subgroup analyses by demographics, infection severity, and
vaccination status that smaller cohorts cannot support.

% =============================================================================
\section{Immune Profiling and T~Cell Research}
\label{sec:registry-immune-profiling}
% =============================================================================

\subsection{CD8+ T~Cell Exhaustion Atlas}
\label{sec:registry-tcell-exhaustion}

\begin{description}
    \item[Principal Investigator:] Multiple groups (NIH intramural)
    \item[Institution:] NIH / multi-centre
    \item[Contact/URL:] \url{https://doi.org/10.1073/pnas.2415119121}
    \item[Funder:] NIH intramural
    \item[Status:] Published 2025
    \item[Phase:] Observational (scRNA-seq atlas)
    \item[Cohort:] 28 ME/CFS patients + 30 sedentary controls; 336,269
        T~lymphoid cells analysed
    \item[Mechanism/Focus:] Single-cell transcriptomics of CD8+ T~cells at
        baseline and after symptom provocation (exercise challenge).
        Identified transcriptional reprogramming toward exhaustion phenotypes,
        with altered effector, memory, and regulatory T~cell subsets.
    \item[Primary Outcomes:] CD8+ T~cells in ME/CFS show exhaustion-associated
        gene expression; exercise provocation amplifies the dysregulation
    \item[Estimated Completion:] Published
    \item[Publication Medium:] \textit{Proceedings of the National Academy
        of Sciences}
    \item[Document Relevance:] Ch.~7 (immune dysfunction), Ch.~6 (energy
        metabolism---immune cell energetics), Ch.~22 (mechanistic studies)
\end{description}

T~cell exhaustion was previously studied mainly in cancer and chronic viral
infections.  Its identification in ME/CFS opens a therapeutic avenue:
checkpoint inhibitor-like approaches that reverse exhaustion might restore
immune surveillance and reduce symptom
burden~\cite{Mayer2025Tcellexhaustion}.  The exercise provocation design is
particularly valuable, as it captures the dynamic immune dysregulation
underlying post-exertional malaise rather than just resting-state
differences.  The finding that monocyte fractions correlate with disease
severity suggests potential for severity stratification.

\subsection{Single-Cell Immune Remodelling in Long COVID--ME/CFS}
\label{sec:registry-elahi-singlecell}

\begin{description}
    \item[Principal Investigator:] Shokrollah Elahi and collaborators
    \item[Institution:] University of Alberta, Canada
    \item[Contact/URL:] Science for ME discussion thread
    \item[Funder:] Institutional / CIHR
    \item[Status:] Published 2026
    \item[Phase:] Observational (scRNA-seq)
    \item[Cohort:] Female long COVID--ME/CFS patients (12~months post-acute
        infection) versus recovered individuals
    \item[Mechanism/Focus:] scRNA-seq of PBMCs examining monocyte, NK~cell,
        and T~cell compartments.  Identified persistent immune remodelling
        12~months after infection in those who developed ME/CFS but not in
        those who recovered.
    \item[Primary Outcomes:] Distinct monocyte polarisation, NK~cell
        dysfunction, and T~cell exhaustion patterns specific to ME/CFS
        development after COVID-19
    \item[Estimated Completion:] Published
    \item[Publication Medium:] \textit{Nature Communications}
    \item[Document Relevance:] Ch.~7 (immune dysfunction), Ch.~14
        \S\ref{sec:cross-disease} (long COVID comparison)
\end{description}

This study complements the NIH T~cell atlas
(Section~\ref{sec:registry-tcell-exhaustion}) by examining the transition
from acute infection to chronic disease.  The comparison of ME/CFS
developers versus recoverers at the same time point (12~months) provides a
natural experiment for identifying immune features that distinguish
pathological from physiological post-infection states~\cite{Elahi2026singlecell}.

\subsection{IgG Immune Complexes Disrupting Cellular Energetics}
\label{sec:registry-igg-complexes}

\begin{description}
    \item[Principal Investigator:] Multiple groups
    \item[Institution:] Multi-centre
    \item[Funder:] Institutional
    \item[Status:] Published 2025
    \item[Phase:] Observational (in vitro)
    \item[Cohort:] ME/CFS patients (including post-COVID ME/CFS) + controls
    \item[Mechanism/Focus:] Isolated IgG complexes from ME/CFS patient
        sera and applied them to healthy cells.  Demonstrated that
        ME/CFS-derived IgG disrupts cellular energetics (mitochondrial
        respiration) and alters inflammatory marker secretion, suggesting
        a ``something in the blood'' pathomechanism.
    \item[Primary Outcomes:] IgG from ME/CFS patients impairs healthy cell
        energy production; effect present in both classical and post-COVID
        ME/CFS
    \item[Estimated Completion:] Published
    \item[Publication Medium:] \textit{Journal of Neuroinflammation}
    \item[Document Relevance:] Ch.~6 (energy metabolism), Ch.~7 (immune
        dysfunction---autoantibodies), Ch.~13 (integrative models)
\end{description}

This is among the strongest ``transfer experiment'' results in ME/CFS:
applying patient-derived IgG to healthy cells recapitulates the energy deficit,
implicating humoral immune factors as causal rather than
epiphenomenal~\cite{IgGcomplexes2025}.  If replicated with purified IgG
subfractions, it could identify specific autoantibody targets and converge
with the daratumumab (Section~\ref{sec:registry-resetme}) and
immunoadsorption (Section~\ref{sec:registry-charite-ia}) therapeutic
approaches.

% =============================================================================
\section{Cardiovascular, Autonomic, and Vascular Research}
\label{sec:registry-cardiovascular}
% =============================================================================

\subsection{MCAM Multi-Site Autonomic Dysfunction Study}
\label{sec:registry-mcam}

\begin{description}
    \item[Principal Investigator:] Elizabeth Unger and collaborators
    \item[Institution:] CDC / multi-site (USA)
    \item[Contact/URL:] \url{https://doi.org/10.3390/jcm14176269}
    \item[Funder:] CDC
    \item[Status:] Published 2025
    \item[Phase:] Observational (multi-site clinical assessment)
    \item[Cohort:] ME/CFS patients from the MCAM (Multi-Site Clinical
        Assessment of ME/CFS) programme
    \item[Mechanism/Focus:] Characterised autonomic dysfunction using
        COMPASS-31 scores and the NASA lean test.  Demonstrated high
        prevalence of autonomic symptoms across ME/CFS cohorts regardless
        of recruitment site.
    \item[Primary Outcomes:] Autonomic dysfunction symptoms are pervasive
        in ME/CFS; COMPASS-31 scores significantly elevated versus controls
    \item[Estimated Completion:] Published
    \item[Publication Medium:] \textit{Journal of Clinical Medicine}
    \item[Document Relevance:] Ch.~10 (cardiovascular), Ch.~8 (neurological
        dysfunction---autonomic)
\end{description}

The MCAM study's multi-site design strengthens generalisability: autonomic
dysfunction is not an artefact of specialist referral patterns but a
consistent feature across diverse clinical
settings~\cite{MCAM2025autonomic}.  The NASA lean test used here is simpler
than head-up tilt, making it a practical screening tool for primary care.

\subsection{Small Fibre Neuropathy in ME/CFS and Post-COVID}
\label{sec:registry-sfn}

\begin{description}
    \item[Principal Investigator:] Naia Azcue and collaborators
    \item[Institution:] Multi-centre (Spain)
    \item[Contact/URL:] \url{https://doi.org/10.1111/ene.70016}
    \item[Funder:] Institutional
    \item[Status:] Published 2025
    \item[Phase:] Observational (case-control)
    \item[Cohort:] ME/CFS patients, post-COVID patients, and healthy controls
    \item[Mechanism/Focus:] Assessed small fibre neuropathy using
        electrochemical skin conductance, nerve reaction times to temperature
        stimuli, and corneal confocal microscopy.  Both ME/CFS and post-COVID
        patients showed significantly higher autonomic and neuropathic symptom
        burden than controls.
    \item[Primary Outcomes:] SFN features present in both ME/CFS and
        post-COVID; corneal nerve fibre density reduced
    \item[Estimated Completion:] Published
    \item[Publication Medium:] \textit{European Journal of Neurology}
    \item[Document Relevance:] Ch.~8 (neurological dysfunction), Ch.~10
        (cardiovascular---perfusion), Ch.~14 \S\ref{sec:cross-disease}
\end{description}

Small fibre neuropathy provides an objective, quantifiable marker for a
symptom---neuropathic pain and dysautonomia---that is otherwise
self-reported~\cite{Azcue2025sfn}.  The corneal confocal microscopy approach
is non-invasive and could serve as a bedside diagnostic.  The shared SFN
phenotype between ME/CFS and post-COVID strengthens the case for common
neurovascular pathology.

\subsection{Microvascular Remodelling and Endothelial Dysfunction}
\label{sec:registry-microvascular}

\begin{description}
    \item[Principal Investigator:] Multiple groups
    \item[Institution:] Multi-centre (Europe)
    \item[Contact/URL:] \url{https://doi.org/10.64898/2026.01.22.26344661}
    \item[Funder:] Institutional
    \item[Status:] Preprint January 2026
    \item[Phase:] Observational (case-control)
    \item[Cohort:] Post-COVID and ME/CFS patients versus healthy controls
    \item[Mechanism/Focus:] Used retinal vessel analysis as a non-invasive
        measure of systemic microvascular health.  Identified reduced venular
        flicker-induced dilation in patients, indicating ongoing endothelial
        dysfunction even in clinically stable individuals.
    \item[Primary Outcomes:] Endothelial dysfunction detectable via retinal
        imaging; reduced venular reactivity in ME/CFS and post-COVID
    \item[Estimated Completion:] Preprint published; peer review pending
    \item[Publication Medium:] medRxiv
    \item[Document Relevance:] Ch.~10 (cardiovascular), Ch.~13 (integrative
        models)
\end{description}

Retinal vessel analysis offers a window into systemic microvascular health
without invasive procedures~\cite{MicrovascularRemodeling2026}.  The
persistent endothelial dysfunction in clinically stable patients suggests
ongoing vascular injury that may underlie fatigue and cognitive symptoms
even when patients report partial recovery.  This complements the microclot
findings~\cite{Nunes2022microclots} and the coagulation proteomics
data~\cite{CoagulationProteomics2024}.

\subsection{POTS Autonomic Profiling Across Aetiologies}
\label{sec:registry-pots-profiling}

\begin{description}
    \item[Principal Investigator:] Multiple groups
    \item[Institution:] Multi-centre
    \item[Status:] Published 2025
    \item[Phase:] Observational (comparative)
    \item[Cohort:] POTS patients from syncope, CFS, and post-COVID-19
        populations
    \item[Mechanism/Focus:] Identified distinct autonomic profiles:
        syncope-related POTS showed baroreceptor dysfunction and sympathetic
        predominance, whereas CFS-related POTS was characterised by
        parasympathetic impairment and impaired short-term baroreflex
        regulation.
    \item[Primary Outcomes:] CFS-related POTS has a distinct autonomic
        signature from syncope-related and post-COVID POTS
    \item[Estimated Completion:] Published
    \item[Document Relevance:] Ch.~10 (cardiovascular), Ch.~8 (neurological
        dysfunction---autonomic)
\end{description}

The distinction between CFS-related and syncope-related POTS autonomic
profiles has therapeutic implications: parasympathetic impairment in
CFS-related POTS suggests vagal tone-targeting interventions (e.g., tVNS,
Section~\ref{sec:registry-tvns}) may be more appropriate than the
beta-blockade used for hyperadrenergic
POTS~\cite{POTSautonomicProfiles2025}.

% =============================================================================
\section{Exercise Physiology and Muscle Biology}
\label{sec:registry-exercise}
% =============================================================================

\subsection{Largest Two-Day CPET Study}
\label{sec:registry-2day-cpet}

\begin{description}
    \item[Principal Investigator:] Multiple groups
    \item[Institution:] Multi-centre (USA)
    \item[Contact/URL:] \url{https://doi.org/10.1186/s12967-024-05410-5}
    \item[Funder:] Institutional
    \item[Status:] Published 2024
    \item[Phase:] Observational (physiological testing)
    \item[Cohort:] Largest 2-day CPET cohort in ME/CFS to date
    \item[Mechanism/Focus:] Two-day cardiopulmonary exercise testing to
        objectively measure post-exertional malaise.  Significant reductions
        in VO$_2$ and workload at the ventilatory anaerobic threshold (VAT)
        on day~2 in ME/CFS versus controls.  Implicates autonomic
        dysregulation of blood flow and oxygen delivery for energy metabolism.
    \item[Primary Outcomes:] Impaired recovery after exertional stressor;
        cardiac, pulmonary, and metabolic contributions to exertion
        intolerance; larger effect sizes at VAT than peak
    \item[Estimated Completion:] Published
    \item[Publication Medium:] \textit{Journal of Translational Medicine}
    \item[Document Relevance:] Ch.~6 (energy metabolism), Ch.~10
        (cardiovascular), Ch.~20 (biomarker research---2-day CPET as
        objective PEM marker)
\end{description}

The 2-day CPET protocol remains the gold standard for objectively
demonstrating PEM.  This largest study to date substantiates that the
day-2 decrement is reproducible across diverse ME/CFS
cohorts~\cite{keller2024cpet}.  The finding that cardiac, pulmonary, and
metabolic factors all contribute to exertion intolerance supports the
multi-system model in Ch.~13 rather than any single-organ explanation.

\subsection{Skeletal Muscle Properties in Long COVID and ME/CFS}
\label{sec:registry-skeletal-muscle}

\begin{description}
    \item[Principal Investigator:] Multiple groups
    \item[Institution:] Multi-centre (Netherlands / Europe)
    \item[Contact/URL:] \url{https://doi.org/10.1101/2025.05.02.25326885}
    \item[Funder:] Institutional
    \item[Status:] Preprint May 2025
    \item[Phase:] Observational (muscle biopsy)
    \item[Cohort:] Long COVID and ME/CFS patients versus controls
    \item[Mechanism/Focus:] Muscle biopsies analysed for mitochondrial
        markers, WASF3 protein, ER stress markers.  Found elevated WASF3 and
        ER stress with decreased mitochondrial Complex~IV subunits.
        Critically, muscle abnormalities differed from those seen in bed rest
        deconditioning.
    \item[Primary Outcomes:] ME/CFS muscle pathology is distinct from
        deconditioning; WASF3 elevation confirmed in skeletal muscle tissue
    \item[Estimated Completion:] Preprint published; peer review pending
    \item[Publication Medium:] medRxiv
    \item[Document Relevance:] Ch.~6 (energy metabolism---WASF3), Ch.~17
        (lifestyle interventions---exercise contraindications)
\end{description}

This study provides crucial evidence against the deconditioning hypothesis:
if ME/CFS fatigue were due to simple disuse, muscle changes should resemble
bed rest patterns~\cite{SkeletalMuscle2025longCOVID}.  The distinct
pathology (WASF3 elevation, mitochondrial complex loss) aligns with the NIH
WASF3 findings~\cite{wang2023wasf3} and supports a disease-specific
metabolic myopathy rather than secondary deconditioning.

% =============================================================================
\section{Microbiome and Gut--Brain Axis}
\label{sec:registry-microbiome}
% =============================================================================

\subsection{JAX AI-Integrated Gut Microbiome Diagnostic}
\label{sec:registry-jax-microbiome}

\begin{description}
    \item[Principal Investigator:] Jackson Laboratory and Duke University
        collaborators
    \item[Institution:] JAX / Duke University (USA)
    \item[Contact/URL:] \url{https://www.jax.org/news-and-insights/2025/july/gut-microbiome-may-predict-invisible-chronic-fatigue-syndrome-and-long-covid}
    \item[Funder:] Institutional / NIH
    \item[Status:] Published July 2025
    \item[Phase:] Diagnostic development (AI-driven)
    \item[Cohort:] 249 individuals
    \item[Mechanism/Focus:] Applied an AI platform to stool, blood, and
        routine lab data to identify disease-specific biomarkers.  Achieved
        90\% accuracy in distinguishing ME/CFS from healthy controls using
        microbiome and metabolite features.
    \item[Primary Outcomes:] 90\% diagnostic accuracy; multimodal data
        (stool + blood + labs) outperforms any single data type
    \item[Estimated Completion:] Published
    \item[Document Relevance:] Ch.~11 (gut microbiome), Ch.~20 (biomarker
        research), Ch.~25 (translational findings)
\end{description}

The 90\% accuracy is the highest reported for any ME/CFS biomarker panel to
date, though it awaits independent validation in a blinded, multi-site
cohort~\cite{JAXmicrobiome2025}.  The multimodal approach (combining stool
microbiome with blood and clinical data) may explain the performance gain
over single-modality studies.  This directly informs the diagnostic algorithm
discussed in Ch.~20.

\subsection{Gut Microbial Composition Differences in CFS}
\label{sec:registry-gut-composition}

\begin{description}
    \item[Principal Investigator:] Multiple groups
    \item[Institution:] Multi-centre
    \item[Contact/URL:] \url{https://doi.org/10.1038/s41598-025-16438-y}
    \item[Status:] Published 2025
    \item[Phase:] Observational (case-control)
    \item[Cohort:] CFS patients versus healthy controls
    \item[Mechanism/Focus:] 16S rRNA sequencing of gut microbiota.  Found
        reduced microbial diversity, depletion of anti-inflammatory taxa
        (\textit{Faecalibacterium prausnitzii}, \textit{Bifidobacterium}),
        and enrichment of pro-inflammatory \textit{Clostridium} species in
        ME/CFS.
    \item[Primary Outcomes:] Significant between-group differences in gut
        microbial composition
    \item[Estimated Completion:] Published
    \item[Publication Medium:] \textit{Scientific Reports}
    \item[Document Relevance:] Ch.~11 (gut microbiome), Ch.~7 (immune
        dysfunction---gut-mediated inflammation)
\end{description}

The depletion of \textit{F.\ prausnitzii} (a major butyrate producer) is
particularly notable because butyrate maintains intestinal barrier integrity.
Its loss could contribute to the intestinal permeability (``leaky gut'')
reported in ME/CFS, creating a pathway from microbiome dysbiosis to systemic
immune activation~\cite{GutMicrobialComposition2025}.

% =============================================================================
\section{Endocrine and Metabolic Profiling}
\label{sec:registry-endocrine}
% =============================================================================

\subsection{Sex-Specific Steroid Hormone Profiling}
\label{sec:registry-steroid-hormones}

\begin{description}
    \item[Principal Investigator:] Multiple groups
    \item[Institution:] Multi-centre
    \item[Contact/URL:] \url{https://doi.org/10.1007/s40618-024-02334-1}
    \item[Funder:] Institutional
    \item[Status:] Published 2024
    \item[Phase:] Observational (cross-sectional)
    \item[Cohort:] ME/CFS patients + healthy controls
    \item[Mechanism/Focus:] UHPLC-MS/MS analysis of circulating steroid
        hormones: mineralocorticoids (aldosterone), glucocorticoids
        (cortisol, corticosterone, 11-deoxycortisol, cortisone), androgens
        (androstenedione, testosterone), and progestins (progesterone,
        17$\alpha$-hydroxyprogesterone).  Found sex-based differences in
        hormone profiles.
    \item[Primary Outcomes:] Sex-specific steroid hormone dysregulation in
        ME/CFS; reinforces need for sex-stratified analyses
    \item[Estimated Completion:] Published
    \item[Publication Medium:] \textit{Journal of Endocrinological
        Investigation}
    \item[Document Relevance:] Ch.~8 (neuroendocrine dysfunction), Ch.~23
        (epidemiology---sex differences)
\end{description}

The sex-based differences in steroid profiles may partly explain the
female predominance in ME/CFS ($\sim$3:1 female-to-male
ratio)~\cite{pipper2024steroid}.  The use of mass spectrometry
(rather than immunoassay) provides superior specificity for closely related
steroid metabolites, reducing cross-reactivity artefacts that have
confounded earlier HPA axis studies.

% =============================================================================
\section{Kynurenine Pathway}
\label{sec:registry-kynurenine}
% =============================================================================

\subsection{OMF Kynurenine Supplementation Trial}
\label{sec:registry-kynurenine-trial}

\begin{description}
    \item[Principal Investigator:] Open Medicine Foundation researchers
    \item[Institution:] OMF-funded site
    \item[Contact/URL:] \url{https://www.omf.ngo/kynurenine-trial-in-me-cfs/}
    \item[Funder:] Open Medicine Foundation
    \item[Status:] Recruiting / active
    \item[Phase:] Interventional (RCT, crossover)
    \item[Cohort:] ME/CFS patients
    \item[Mechanism/Focus:] Randomised to kynurenine or placebo for 3~months
        followed by crossover.  Tests whether exogenous kynurenine
        supplementation can correct tryptophan metabolism disturbances
        reported in ME/CFS.  The kynurenine pathway mediates immune response
        and neuroinflammation through metabolites like quinolinic acid and
        kynurenic acid.
    \item[Primary Outcomes:] Fatigue, cognitive function, immune markers;
        kynurenine/tryptophan ratio changes
    \item[Estimated Completion:] TBD
    \item[Document Relevance:] Ch.~6 (energy metabolism---NAD+ synthesis
        via kynurenine pathway), Ch.~8 (neurological dysfunction---
        neuroinflammation), Ch.~16 (supplements)
\end{description}

The kynurenine pathway sits at the intersection of immune activation and
energy metabolism: IDO-driven tryptophan diversion reduces serotonin
availability while generating both neuroprotective (kynurenic acid) and
neurotoxic (quinolinic acid) metabolites~\cite{OMFkynureninetrial}.
If kynurenine supplementation improves symptoms, it would confirm that
tryptophan diversion is a treatable metabolic bottleneck in ME/CFS.

% =============================================================================
\section{Comorbid Conditions: EDS--MCAS--ME/CFS Triad}
\label{sec:registry-comorbid-triad}
% =============================================================================

\subsection{EDS--MCAS Genetic and Pathophysiological Basis}
\label{sec:registry-eds-mcas}

\begin{description}
    \item[Principal Investigator:] Multiple groups
    \item[Institution:] Multi-centre (international)
    \item[Contact/URL:] PMID 41554340 (2026 review); PMID 39451569 (2024
        genetics)
    \item[Funder:] Institutional
    \item[Status:] Published 2024--2026
    \item[Phase:] Review and genetic analysis
    \item[Cohort:] hEDS patients + controls (genetics: $n \approx 250$)
    \item[Mechanism/Focus:] Two convergent lines of research: (1) a 2026
        review examining pathophysiological and clinical links between EDS
        and MCAS, highlighting bidirectional fibroblast--mast cell
        interactions in the extracellular matrix; (2) a 2024 genetic study
        identifying variants in MT-CYB, HTT, MUC3A, HLA-B, and HLA-DRB1
        implicated in hEDS and mast cell hypersensitivity (72.2\% of hEDS
        cohort vs 14.2\% controls).
    \item[Primary Outcomes:] Genetic basis for the EDS--MCAS overlap;
        pathophysiological framework for the clinical triad
    \item[Estimated Completion:] Published
    \item[Document Relevance:] Ch.~7 (immune dysfunction---mast cells),
        Ch.~10 (cardiovascular---connective tissue), Ch.~14
        \S\ref{sec:cross-disease} (comorbidity patterns)
\end{description}

The EDS--MCAS--ME/CFS triad is one of the most clinically recognisable
comorbidity patterns, yet its biological basis has been poorly
understood~\cite{EDSMCAS2026review, hEDSgenetics2024MCAS}.  The
fibroblast--mast cell interaction framework provides a mechanistic link:
abnormal collagen (EDS) alters the extracellular matrix, which modulates
mast cell behaviour (MCAS), which drives systemic immune activation
(ME/CFS).  The HLA associations suggest an autoimmune component that
connects to the broader autoantibody research in
Sections~\ref{sec:registry-charite-ia}--\ref{sec:registry-resetme}.

% =============================================================================
\section{Conferences and Funding Initiatives}
\label{sec:registry-conferences}
% =============================================================================

\subsection{International ME/CFS Conference 2026 (Berlin)}
\label{sec:registry-berlin-2026}

\begin{description}
    \item[Organiser:] ME/CFS Research Foundation
    \item[Location:] Berlin, Germany
    \item[Dates:] 7--8 May 2026
    \item[Contact/URL:] \url{https://mecfs-research.org/en/news-conference2026/}
    \item[Focus:] Day~1: disease mechanisms (PEM, pain, cognitive dysfunction,
        circulatory disorders, dysautonomia).  Day~2: clinical trial results
        and treatment data.  Aimed at the medical and research community.
    \item[Document Relevance:] Multiple chapters; unpublished data presented
        here may update several registry entries
\end{description}

This conference is the primary venue for new ME/CFS data presentation in
2026~\cite{BerlinConference2026}.  Several registry entries (BC007, HBOT,
LDN trials) may present interim or final results.  The registry should be
updated after the conference.

\subsection{ME/CFS Research Foundation 2026 Funding Strategy}
\label{sec:registry-mecfs-foundation-2026}

\begin{description}
    \item[Organiser:] ME/CFS Research Foundation (Germany)
    \item[Contact/URL:] \url{https://mecfs-research.org/en/fundingstrategy2026/}
    \item[Focus:] 2026 funding priorities: autoimmunity, immune
        dysregulation, metabolic disorders, inflammation, vascular
        dysfunction, and nervous system dysfunction as causal factors in
        ME/CFS.  Supports projects addressing central symptoms (PEM, pain,
        cognitive dysfunction, circulatory disorders, dysautonomia).
    \item[Document Relevance:] Signals research direction for the field;
        funded projects will appear as future registry entries
\end{description}

% =============================================================================
\section{Digital Health and Wearable Technology}
\label{sec:registry-digital}
% =============================================================================

\subsection{Long Covid Wearable Device Study}
\label{sec:registry-wearable}

\begin{description}
    \item[Principal Investigator:] Julia Moore Vogel, PhD
    \item[Institution:] Scripps Research Institute, United States
    \item[Contact/URL:] \url{https://solvecfs.org/exploring-current-me-cfs-and-long-covid-clinical-trials-on-national-clinical-trials-day/}
    \item[Funder:] Institutional
    \item[Status:] Recruiting
    \item[Phase:] Observational, digital health
    \item[Cohort:] Adults with Long Covid, ME/CFS, or POTS
    \item[Mechanism/Focus:] Evaluates whether consumer wearable devices
        (e.g., Fitbit, Garmin) can support activity pacing, symptom
        management, and overall quality of life.  Collects continuous
        physiological data (heart rate, activity, sleep) to quantify
        PEM triggers.
    \item[Primary Outcomes:] Pacing adherence; symptom management
        effectiveness; quality of life; PEM episode frequency
    \item[Estimated Completion:] Not yet announced
    \item[Timeline:] Recruitment ongoing as of 2026
    \item[Publication Medium:] Not yet announced
    \item[Document Relevance:] Ch.~17 (lifestyle interventions---pacing),
        Ch.~10 (cardiovascular/autonomic monitoring)
\end{description}

While not a biomarker or therapeutic study, this research addresses a critical
practical question for ME/CFS management.  Activity pacing is the primary
non-pharmacological intervention (Chapter~17), yet objective tools to support
pacing remain poorly validated.  If wearable-derived metrics (heart rate
variability, activity thresholds) can reliably predict PEM onset, this could
transform day-to-day management for patients.

% =============================================================================
\section{Results Expected in 2026--2027}
\label{sec:registry-upcoming}
% =============================================================================

Several studies in this registry are expected to report results in the near
term.  Table~\ref{tab:registry-upcoming} highlights those with the highest
potential impact on ME/CFS understanding and treatment.

\begin{table}[htbp]
\centering
\small
\caption{Studies with results expected in 2026--2027, ranked by anticipated
impact.}
\label{tab:registry-upcoming}
\begin{tabular}{@{}p{3.0cm}p{2.0cm}p{7.3cm}@{}}
\toprule
\textbf{Study} & \textbf{Expected} & \textbf{Why it matters} \\
\midrule
DecodeME & Mid--late 2026 &
    First adequately powered GWAS ($n>20{,}000$); could identify
    causal genetic loci and enable Mendelian randomisation \\
BC007 Phase II & 2026 &
    Largest RCT ($n=114$) testing autoantibody-targeted therapy;
    validates or refutes the functional GPCR autoantibody hypothesis \\
RECOVER-TLC & End of 2026 &
    First-round results (baricitinib, structured pacing); largest
    federally funded treatment trials relevant to ME/CFS \\
BioQuest & 2026--2027 &
    Largest ME/CFS biomarker study ($n=1{,}000$); could yield first
    validated blood-based diagnostic panel with disease controls \\
BioMapAI & 2026--2027 &
    Longitudinal multi-omics follow-up; will distinguish state vs
    trait markers over 4-year disease trajectory \\
ADDRESS-LC & H1 2026 &
    First CNS-penetrant NF-$\kappa$B modulator for Long Covid brain fog;
    cross-disease (Alzheimer's, Parkinson's) data available \\
Cornell NAC-GSH & $\sim$2026 &
    First dose-response study of NAC on brain glutathione (MRS);
    could validate oxidative stress as druggable CNS target \\
NSU Probiotics & $\sim$2026 &
    First ME/CFS-specific probiotic RCT with IBS stratification;
    tests gut--brain axis hypothesis directly \\
OMF Kynurenine & $\sim$2026 &
    First RCT testing kynurenine supplementation; probes
    tryptophan--serotonin--neuroinflammation axis \\
Berlin Conference & May 2026 &
    Major venue for unpublished data; BC007, HBOT, and LDN results
    may be presented \\
Nacul LDN RCT & Early 2026 &
    First double-blind ME/CFS-specific LDN trial with inflammatory
    markers and QoL endpoints \\
\bottomrule
\end{tabular}
\end{table}

\textbf{Recently published results to monitor for follow-up studies.}
Multiple registry entries have reported results that may generate
replication or extension studies in 2026--2027:
\begin{itemize}
    \item \textbf{RESTORE-ME} (oxaloacetate): Positive RCT ($p=0.0039$ for
        fatigue reduction); independent replication needed.
    \item \textbf{STOP-PASC / PAX~LC} (Paxlovid): Negative results challenge
        the viral persistence hypothesis; redirects research toward
        post-infectious immune and metabolic mechanisms.
    \item \textbf{Stellate ganglion block}: Promising pilot ($n=10$) for
        autonomic reset; larger sham-controlled RCT likely in planning.
    \item \textbf{HBOT pilot}: Thalamic hyperconnectivity normalised in 30
        ME/CFS patients; sham-controlled RCT needed.
    \item \textbf{Cornell cfRNA}: Liquid biopsy classifier (AUC 0.81);
        multi-site validation with disease controls required.
    \item \textbf{Palomo herpesvirus reactivation}: Multi-virus co-reactivation
        in 72.5\% of patients; anti-dUTPase IgG as potential stratification
        marker.
    \item \textbf{JAX microbiome AI}: 90\% diagnostic accuracy from multimodal
        gut/blood data; blinded independent replication needed.
    \item \textbf{Post-IM 7-year cohort}: 4\% locked-in ME/CFS rate at
        24~months; risk factor analysis pending.
\end{itemize}

% =============================================================================
\section{Registry Summary}
\label{sec:registry-summary}
% =============================================================================

Table~\ref{tab:registry-summary} provides a compact overview of all tracked
studies.

\begin{table}[htbp]
\centering
\small
\caption{Summary of ongoing and planned ME/CFS research studies.}
\label{tab:registry-summary}
\begin{tabular}{@{}p{3.0cm}p{2.6cm}p{1.6cm}p{1.3cm}p{4.8cm}@{}}
\toprule
\textbf{Study} & \textbf{Institution} & \textbf{Status} & \textbf{Cohort} & \textbf{Focus} \\
\midrule
Rosetta Stone & Imperial College London & Active & $n=500+$ & Multi-omics: ME/CFS vs Long Covid \\
BioQuest & Uppsala / OMF & Analysis & $n=1{,}000$ & Multi-omics biomarker panel \\
DecodeME & Univ.\ Edinburgh & Analysis & $n>20{,}000$ & GWAS for ME/CFS \\
EpiSwitch CFS & Oxford BioDynamics & Validation & $n=108^*$ & Epigenetic diagnostic test \\
BioMapAI & JAX / Duke & Ongoing & $n=249$ & AI-integrated multi-omics \\
MEA PhD Project & Rosalind Franklin Inst. & Active & TBD & Metabolites \& infection markers \\
Charit\'{e} IA & Charit\'{e} Berlin & Active & TBD & Immunoadsorption (autoantibodies) \\
BC007 Phase II & Berlin Cures & Analysis & $n=143$ & Autoantibody neutralisation \\
NCNED TRPM3 & Griffith Univ. & Active & TBD & Ion channels + LDN trials \\
Charit\'{e} fMRI & Charit\'{e} Berlin & Funded & TBD & Neuroimaging biomarkers \\
RECOVER-TLC & NIH multi-site & Active & $n=600+$ & Baricitinib, LDN, GLP-1, pacing \\
LIFT Trial & Brigham \& Women's & Recruiting & TBD & LDN + pyridostigmine \\
Rapamycin & Multi-site (US) & Active & TBD & mTOR inhibition \\
Lumbrokinase & Mount Sinai & Recruiting & TBD & Fibrinolytic (microclots) \\
RESTORE-ME & Terra Biological & Published & $n=82$ & Oxaloacetate (TCA cycle) \\
Stellate Ganglion & Multi-site (US) & Pilot & $n=10$ & Autonomic reset (SGB) \\
tVNS & Multi-site & Recruiting & TBD & Vagus nerve stimulation \\
Masitinib & AB Science & Phase 2 & TBD & Mast cell (MCAS) \\
STOP-PASC/PAX LC & NIH multi-site & Published & TBD & Paxlovid (negative result) \\
Hydrogen Water & Stony Brook & Recruiting & TBD & Antioxidant dosing \\
Wearable Study & Scripps Research & Recruiting & TBD & Pacing via wearables \\
\midrule
\multicolumn{5}{@{}l}{\textit{Added March 2026}} \\
\midrule
MEA--UCL Biomarker & UCL / MEA & Active & Biobank & Blood biomarkers (neuro) \\
ResetME & Haukeland / Oslo & Recruiting & $n=66$ & Daratumumab (anti-CD38) \\
UAB LDN & U.\ Alabama (Younger) & Recruiting & TBD & LDN dose-finding \\
Cornell NAC-GSH & Weill Cornell & Recruiting & $n=60$ & Brain glutathione (MRS) \\
NSU Probiotics & Nova Southeastern & Recruiting & $n=100$ & Microbiome / gut--brain axis \\
ADDRESS-LC & BioVie / DoD & Recruiting & $n=200$ & Neuroinflammation (brain fog) \\
LC-REVITALIZE & Western Univ.\ / SILC & Recruiting & $n=348$ & JAK1 + anti-fibrotic \\
CHIIME & UCSF / PolyBio & Recruiting & TBD & Viral tissue persistence \\
\midrule
\multicolumn{5}{@{}l}{\textit{Added March 2026 (second expansion)}} \\
\midrule
PrecisionLife & PrecisionLife / DecodeME & Published & $n>20{,}000$ & Combinatorial genomics (250+ genes) \\
Cornell cfRNA & Cornell Univ. & Published & $n=168$ & Cell-free RNA liquid biopsy \\
CRM Multi-modal & Multi-centre (US) & Published & varied & Energy + immune + vascular \\
HBOT Pilot & Multi-site & Published & $n=30$ & Hyperbaric oxygen therapy \\
Palomo Herpesvirus & Multi-centre & Published & varied & EBV/HHV-6/VZV co-reactivation \\
Post-IM 7yr & Northwestern & Published & $n=4{,}501$ & Post-mononucleosis ME/CFS cohort \\
RECOVER Incidence & NIH multi-site & Published & large & Post-COVID ME/CFS incidence \\
CD8+ T~cell Atlas & NIH intramural & Published & $n=58$ & scRNA-seq T~cell exhaustion \\
Elahi scRNA-seq & U.\ Alberta & Published & varied & Immune remodelling (long COVID) \\
IgG Complexes & Multi-centre & Published & varied & Transfer experiment (energetics) \\
MCAM Autonomic & CDC multi-site & Published & varied & COMPASS-31 + lean test \\
SFN (Azcue) & Multi-centre (Spain) & Published & varied & Small fibre neuropathy \\
Microvascular & Multi-centre (EU) & Preprint & varied & Retinal vessel / endothelial \\
POTS Profiling & Multi-centre & Published & varied & Autonomic subtyping \\
2-Day CPET & Multi-centre (US) & Published & largest & PEM objective measurement \\
Muscle Biopsy & Multi-centre (EU) & Preprint & varied & WASF3 / ER stress in muscle \\
JAX Microbiome AI & JAX / Duke & Published & $n=249$ & Gut diagnostic (90\% accuracy) \\
Gut Composition & Multi-centre & Published & varied & Microbial dysbiosis \\
Steroid Hormones & Multi-centre & Published & varied & Sex-specific HPA profiling \\
OMF Kynurenine & OMF-funded & Recruiting & TBD & Tryptophan metabolism RCT \\
EDS--MCAS & Multi-centre & Published & $n\approx250$ & Triad genetics/pathophysiology \\
Nacul LDN & UBC / BC Women's & Active & TBD & LDN double-blind RCT \\
Rituximab Japan & Japan & Active & $n=30$ & Phase 2 rituximab + imaging \\
Berlin Conf.\ 2026 & ME/CFS Res.\ Found. & May 2026 & --- & Conference (mechanisms + trials) \\
\bottomrule
\multicolumn{5}{@{}l}{\footnotesize $^*$Proof-of-concept; larger validation cohorts planned.}
\end{tabular}
\end{table}
